% =====================================================================
% TÜRKÇE İNCELEME KOPYASI (gönderim sürümü değildir; SAUCIS yalnızca
% İngilizce kabul eder). Kaynak: ../supplementary/supplementary.tex
% "Boyuttan Bağımsız V2G Şarj Kontrolü için Olurluluğu Koruyan Güvenlik
%  Katmanına Sahip Dikkat Tabanlı Güvenli Pekiştirmeli Öğrenme" için ek
%  materyal (çalışma başlığı, DD30; nihai başlığı yazar belirler)
% Derleme: latexmk -pdf supplementary_tr.tex (ana metne yapılan çapraz
% göndermelerin çözülmesi için önce ../manuscript_tr/main_tr.tex bir kez
% derlenmelidir)
% =====================================================================
\newif\ifsaucis
\IfFileExists{saucis.cls}{\saucistrue\documentclass[final,9pt,onecolumn]{saucis}}{\saucisfalse\documentclass[10pt]{article}}
\newif\ifanonymous
\anonymousfalse

% Türkçe inceleme kopyasının önsözü (manuscript_tr/main_tr.tex); İngilizce
% ana metin ve ek materyalle ortak önsözün uyarlanmış kopyasıdır.
% ---------------- compatibility preamble ----------------------------
\makeatletter
\ifsaucis\else
  \usepackage[a4paper,margin=2.5cm]{geometry}
  \usepackage[T1]{fontenc}
  \IfFileExists{lmodern.sty}{\usepackage{lmodern}}{}
  \setlength{\parindent}{0pt}          % SAUCIS: no paragraph indentation
  \setlength{\parskip}{0.5\baselineskip}
\fi
% Türkçe dil desteği (kısaltmalar kapalı: matematik ve tablolarda = ve : bozulmasın)
\usepackage[utf8]{inputenc}
\@ifpackageloaded{fontenc}{}{\usepackage[T1]{fontenc}}
\usepackage[turkish,shorthands=off]{babel}
\@ifpackageloaded{amsmath}{}{\usepackage{amsmath}}
\@ifpackageloaded{amssymb}{}{\usepackage{amssymb}}
\@ifpackageloaded{amsthm}{}{%
  \@ifundefined{proof}{}{\let\proof\relax\let\endproof\relax}%
  \usepackage{amsthm}}
\@ifpackageloaded{graphicx}{}{\usepackage{graphicx}}
\@ifpackageloaded{booktabs}{}{\usepackage{booktabs}}
\@ifpackageloaded{longtable}{}{\usepackage{longtable}}
\@ifpackageloaded{xcolor}{}{\usepackage{xcolor}}
\@ifpackageloaded{array}{}{\usepackage{array}}
\@ifpackageloaded{cite}{}{\usepackage{cite}}
\@ifpackageloaded{algorithm}{}{\usepackage{algorithm}}
\@ifpackageloaded{algpseudocode}{}{\@ifpackageloaded{algorithmic}{}{\usepackage{algpseudocode}}}
\@ifpackageloaded{caption}{}{\usepackage{caption}}
\IfFileExists{adjustbox.sty}{\@ifpackageloaded{adjustbox}{}{\usepackage{adjustbox}}}{}% generated result tables use adjustbox
\captionsetup{font=small,labelfont=bf}
\@ifundefined{floatname}{}{\floatname{algorithm}{Algoritma}}
\@ifpackageloaded{eurosym}{}{\IfFileExists{eurosym.sty}{\usepackage[gen]{eurosym}}{\providecommand{\euro}{\texteuro}}}
% cross-references between main text and supplementary material
\@ifpackageloaded{hyperref}{\usepackage{xr}}{\usepackage{xr-hyper}\usepackage[hidelinks]{hyperref}}
\makeatother

% Unicode symbols that generated tables may contain (pdflatex)
\makeatletter
\@ifundefined{DeclareUnicodeCharacter}{}{%
\DeclareUnicodeCharacter{03BB}{\ensuremath{\lambda}}\DeclareUnicodeCharacter{03A3}{\ensuremath{\Sigma}}%
\DeclareUnicodeCharacter{03C3}{\ensuremath{\sigma}}\DeclareUnicodeCharacter{03BC}{\ensuremath{\mu}}%
\DeclareUnicodeCharacter{03C1}{\ensuremath{\rho}}\DeclareUnicodeCharacter{03B1}{\ensuremath{\alpha}}%
\DeclareUnicodeCharacter{03B2}{\ensuremath{\beta}}\DeclareUnicodeCharacter{0394}{\ensuremath{\Delta}}%
\DeclareUnicodeCharacter{03B7}{\ensuremath{\eta}}\DeclareUnicodeCharacter{03C4}{\ensuremath{\tau}}%
\DeclareUnicodeCharacter{2264}{\ensuremath{\le}}\DeclareUnicodeCharacter{2265}{\ensuremath{\ge}}%
\DeclareUnicodeCharacter{2212}{\ensuremath{-}}\DeclareUnicodeCharacter{00B7}{\ensuremath{\cdot}}%
\DeclareUnicodeCharacter{00D7}{\ensuremath{\times}}\DeclareUnicodeCharacter{2248}{\ensuremath{\approx}}%
\DeclareUnicodeCharacter{2020}{\ensuremath{\dagger}}}
\makeatother
\setlength{\emergencystretch}{2em}
% anonymous builds: no creation date or trailer id in the PDF metadata
\ifanonymous\ifdefined\pdfinfoomitdate\pdfinfoomitdate=1\fi\ifdefined\pdftrailerid\pdftrailerid{}\fi\fi

% ---------------- theorem environments ------------------------------
\theoremstyle{plain}
\makeatletter
\@ifundefined{theorem}{\newtheorem{theorem}{Teorem}}{}
\@ifundefined{lemma}{\newtheorem{lemma}{Lemma}}{}
\@ifundefined{proposition}{\newtheorem{proposition}{Önerme}}{}
\@ifundefined{corollary}{\newtheorem{corollary}{Sonuç}}{}
\theoremstyle{definition}
\@ifundefined{assumption}{\newtheorem*{assumption}{Varsayım}}{}
\@ifundefined{remark}{\newtheorem{remark}{Not}
\@ifundefined{definition}{\newtheorem{definition}{Tanım}}{}}{}
\makeatother

% Türkçe başlık adları (babel'in varsayılanları: Özet, Kanıt)
\addto\captionsturkish{\renewcommand{\abstractname}{Öz}\renewcommand{\proofname}{İspat}%
  \renewcommand{\figurename}{Şekil}\renewcommand{\tablename}{Tablo}\renewcommand{\refname}{Kaynaklar}}

% ---------------- project macros ------------------------------------
\newcommand{\RES}[1]{\textcolor{red}{[SONUÇ: #1]}}
\newcommand{\TODO}[1]{\textcolor{blue}{[YAPILACAK: #1]}}
\newcommand{\method}{SA-Attn}          % method name; must match analysis/scripts/names.py
\newcommand{\Eq}[1]{Denklem~\ref{#1}}
\newcommand{\Eqs}[2]{Denklem~\ref{#1}--\ref{#2}}
\newcommand{\Fig}[1]{Şekil~\ref{#1}}
\newcommand{\Tab}[1]{Tablo~\ref{#1}}
\newcommand{\Sec}[1]{Bölüm~\ref{#1}}
\newcommand{\Alg}[1]{Algoritma~\ref{#1}}
\newcommand{\R}{\mathbb R}
\newcommand{\clip}{\operatorname{clip}}
\newcommand{\Ctil}{\tilde C}
\newcommand{\Crev}{C_{\mathrm{rev}}}
\newcommand{\Bh}{\hat B}
\newcommand{\Pnom}{\bar P^{\mathrm{nom}}}
\newcommand{\Pdis}{P^{\mathrm{dis}}}
\newcommand{\Ereq}{E^{\mathrm{req}}}
\newcommand{\pmn}{p^{\min}}
\newcommand{\sm}{\mathrm}

% figure from file, or a framed placeholder if the file is missing
\newcommand{\figorbox}[3]{%
  \IfFileExists{#1}{\includegraphics[width=#2]{#1}}{%
    \fbox{\parbox[c][#3][c]{0.9\linewidth}{\centering\small
      Şekil dosyası bulunamadı: \texttt{\detokenize{#1}}}}}}
% first existing file of three candidates (analysis output, analysis source, local copy)
\newcommand{\figchain}[4]{%
  \IfFileExists{#2}{\includegraphics[width=#1]{#2}}{%
  \IfFileExists{#3}{\includegraphics[width=#1]{#3}}{\figorbox{#4}{#1}{5cm}}}}
% Generated results. \analysisdir points to the analysis output tree of the
% planned paper/ + analysis/ layout; a build may override it before \input.
\providecommand{\analysisdir}{../../analysis}
% generated result table: analysis output, local copy in tables/, or placeholder
\newcommand{\resulttable}[2]{\IfFileExists{#1}{\input{#1}}{#2}}
\newcommand{\restable}[2]{%
  \IfFileExists{\analysisdir/tables/#1.tex}{\input{\analysisdir/tables/#1.tex}}{%
  \IfFileExists{tables/#1.tex}{\input{tables/#1.tex}}{#2}}}
% generated result figure: analysis output, local copy in figures/, or a red placeholder box
\newcommand{\resultfig}[2]{%
  \IfFileExists{\analysisdir/figures/#1}{\includegraphics[width=0.85\linewidth]{\analysisdir/figures/#1}}{%
  \IfFileExists{figures/#1}{\includegraphics[width=0.85\linewidth]{figures/#1}}{%
    \fbox{\parbox[c][5cm][c]{0.85\linewidth}{\centering\small
      \textcolor{red}{SONUÇ ŞEKLİ YER TUTUCUSU}\\[2mm]
      beklenen dosya: \texttt{\detokenize{#1}} (analysis/figures/)\\[1mm] #2}}}}}

\graphicspath{{figures/}}
\externaldocument[M-]{../manuscript_tr/main_tr}

% ek materyal numaralandırması: Bölüm S1, Lemma S1, Tablo S1, Şekil S1, (S1)
\renewcommand{\thesection}{S\arabic{section}}
\renewcommand{\thetable}{S\arabic{table}}
\renewcommand{\thefigure}{S\arabic{figure}}
\renewcommand{\theequation}{S\arabic{equation}}
\renewcommand{\thelemma}{S\arabic{lemma}}
\renewcommand{\theproposition}{S\arabic{proposition}}
\renewcommand{\thealgorithm}{S\arabic{algorithm}}
% algoritma girdisi etiketi (denetim sözcükleri sözde kod sözdizimi olarak İngilizce bırakılmıştır)
\algrenewcommand\algorithmicrequire{\textbf{Girdi:}}
% analiz betiklerinin üreteceği tablo için yer tutucu
\newcommand{\placeholdertable}[3]{%
\begin{table}[t]\centering
\caption{#2 (analiz betikleri çalıştırılana kadar yer tutucudur).}\label{#1}% beklenen dosya: analysis/tables/#3
\footnotesize \RES{\texttt{analysis/scripts/make\_all.py} tarafından üretilen tablo}
\end{table}}

\title{Ek Materyal\\[1mm]\large Boyuttan Bağımsız V2G Şarj Kontrolü için Olurluluğu Koruyan Güvenlik Katmanına Sahip Dikkat Tabanlı Güvenli Pekiştirmeli Öğrenme}
\ifanonymous
  \author{Anonim yazar(lar)}
\else
  \author{Yazar Ad(lar)ı\\ \small \TODO{Kurum(lar)}}
\fi
\date{}

\begin{document}
\maketitle

\begin{center}
\fbox{\parbox{0.92\linewidth}{\small\itshape
Bu belge yazarın incelemesi için hazırlanmış Türkçe çeviridir; gönderim sürümü İngilizcedir. Şekillerin içindeki metinler İngilizce bırakılmıştır.}}
\end{center}

Bu belge ana metni tamamlamaktadır.
Bölüm~\ref{app:proofs} biçimsel sonuçları ispatlamakta, Bölüm~\ref{sec:ev2gym} modellemenin EV2Gym'de nasıl yapıldığını belgelemekte ve Bölüm~\ref{sec:pipeline} talep işleme hattını ve tam bilgili optimumları vermektedir.
Karşılaştırılan yöntemler, eğitim, ayarlama ve istatistikler Bölüm~\ref{sec:training-supp} içinde ele alınmakta; Bölüm~\ref{sec:results-supp} ise ek bulguları bir araya getirmektedir.
Ana metnin denklemlerine, lemmalarına ve şekillerine ana metindeki numaralarıyla atıf yapılmakta; bu belgedeki ögeler ise S önekini taşımaktadır.
Kısaltmalar ana metni izlemektedir: bir şarj noktası işletmecisi (CPO) istasyonunda araçtan şebekeye (V2G) hizmetleri ve fotovoltaik (PV) üretim içeren elektrikli araç (EV) şarjı.
Yöntem, yakınsal politika optimizasyonu (PPO) ile eğitilen ve çok katmanlı algılayıcı (MLP), ekonomik model öngörülü kontrol (eMPC) ve heterojen çizge politikası (EV-GNN) ile karşılaştırılan güvenli dikkat tabanlı denetleyicidir (\method{}).
Diğer kısaltmalar doğrusal program (LP), Karush--Kuhn--Tucker (KKT) koşulları, en az esneklik payı önce (LLF) çizelgelemesi ve konformal tahmindir (CP).
\Tab{tab:nomenclature} başlıca sembolleri listelemektedir.

% ---------------------------------------------------------------------
% Semboller (hakem notu m4)
% ---------------------------------------------------------------------
\begin{table}[t]
\centering
\caption{Başlıca semboller. Birimler köşeli parantez içinde verilmiştir; EV $i$'ye ait büyüklükler şarj ünitesinin indisini taşır.}
\label{tab:nomenclature}
\footnotesize
\setlength{\tabcolsep}{3pt}
\begin{tabular}{@{}>{\raggedright}p{2.3cm}>{\raggedright}p{5.25cm}>{\raggedright}p{2.3cm}>{\raggedright\arraybackslash}p{5.25cm}@{}}
\toprule
Sembol & Anlamı & Sembol & Anlamı \\
\midrule
$N$, $\mathcal N(t)$, $n(t)$ & şarj üniteleri, dolu şarj üniteleri, bağlı EV'ler & $\Delta t$, $T$ & adım uzunluğu [h], ufuk [adım] \\
$t_i^{\sm{arr}}$, $d_i$ & varış adımı, son kontrol edilen adım & $\tau_i(t)$ & kalan eylemler, $d_i-t+1$ \\
$E_i(t)$, $E_i^{\sm{cap}}$, $E_i^{\min}$ & batarya enerjisi, kapasite, minimum [kWh] & $E_i^{\sm{des}}$, $\Ereq_i$ & istenen ve kabul edilen enerji [kWh] \\
$p_i$, $\tilde p_i$ & komut edilen ve raporlanan güç [kW] & $g_i$, $\eta_i^{\sm c}$, $\eta_i^{\sm d}$ & enerji eşlemesi ve verimler \\
$\bar P_i(E)$, $\Pnom_i$ & en yüksek şarj gücü, nominal değeri [kW] & $\pmn_i(E)$, $\Pdis_i$ & fiziksel alt limit, deşarj gücü [kW] \\
$P_{\sm{tr}}$, $\rho_{\sm{tr}}$ & trafo anma gücü, şarj ünitesi başına anma gücü [kW] & $P^{\sm{inf}}$, $P^{\sm{pv}}$ & esnek olmayan yük, PV üretimi [kW] \\
$C$, $\Crev$, $\Ctil$ & ileri yön, ters yön ve sıkılaştırılmış ileri yön limiti [kW] & $\varepsilon_q$ & eyleyici payı, 0.008 kW \\
$c^{\sm{ch}}$, $c^{\sm{dis}}$, $c^{\sm{thr}}$ & şarj, deşarj ve enerji akışı (throughput) fiyatı [\texteuro/kWh] & $\varepsilon_{\sm{sat}}$ & doyma eşiği, 0.001 kW \\
$\varphi_i$, $\Psi_{i,k}$, $R_i$ & tek adımlık tam güç eşlemesi, $k$ katlı eşleme, erişilebilir enerji & $B_i(k)$, $\Bh_i(k)$ & kesin ve sertifikalı eşikler [kWh] \\
$l_i$, $u_i$ & güvenlik katmanının alt ve üst sınırı [kW] & $\chi_i$ & esneklik payı [adım] \\
$x_i$, $a_i$, $P^{\sm{ref}}$ & ham eylem [kW], ham örnek, ölçek 11 kW & $\mathcal K(t)$, $\Pi_t$ & olurlu küme, izdüşüm \\
$\lambda^{*}$, $\mathcal F$ & ortak kaydırma [kW], serbest küme & $b$ & yedek kural bütçesi [kW] \\
$\varsigma$, $\varsigma^{\sm{tr}}$, $m$, $\kappa$ & şarj durumu, şarj gücü azalmasının başlangıcı, çarpan, üs & $z_0(m)$, $\kappa^{*}(m)$ & kesin teslim sabitleri \\
$\varrho_i^{\sm{in}}$, $\varrho_i^{\sm{out}}$ & şarj gücü azalmasının iç ve dış eğimi [1/h] & $G_i^{\max}$ & yuvarlama payı [kWh] \\
$J^{\pi}$, $J^{\sm{pen}}$ & kâr, cezalı kâr [\texteuro] & $J^{\sm{pen,LP}}$, $J^{\sm{LP,ceil}}$ & yumuşak ve katı optimum [\texteuro] \\
$w_E$, $w_P$, $w$ & değerlendirme ve eğitim ceza ağırlıkları [\texteuro/kWh] & $\Delta E_i^{\sm{sh}}$ & ayrılıştaki eksik enerji [kWh] \\
$\xi_i$, $\xi_0$ & şarj ünitesi ve küresel belirteç & $d_m$, $n_L$, $n_H$ & genişlik, katman, başlık sayısı \\
$\theta$, $\omega$ & aktör ve kritik parametreleri & $\beta$ & düzenlileştirici ağırlığı \\
$\gamma$, $\lambda_{\sm{GAE}}$, $\epsilon_{\sm{clip}}$ & indirim faktörü, GAE parametresi, kırpma aralığı & $n_{\sm{mb}}$, $c_{\sm{vf}}$, $c_{\sm{ent}}$ & mini yığın boyutu, kayıp katsayıları \\
$\Lambda_E$, $\Lambda_P$, $\eta_\Lambda$ & Lagrange çarpanları, dual adım büyüklüğü & $\alpha$, $n_{\sm{cal}}$ & kapsama dışı kalma düzeyi, kalibrasyon boyutu \\
$\hat q$, $e_j$, $\hat Q_{1-\alpha}$ & kantil tahmincisi, skor, konformal kantil & $\hat d_i^{\sm{lo}}$, $\tau_i^{\sm{lo}}$ & ayrılış sınırı, kalan adımlar \\
\bottomrule
\end{tabular}
\end{table}

% ---------------------------------------------------------------------
\section{İspatlar}\label{app:proofs}
% ---------------------------------------------------------------------

Bu bölüm ana metindeki biçimsel sonuçları ispatlamakta ve ana metnin atıf yaptığı yardımcı lemmaları ifade etmektedir.
Tüm ifadeler her EV için ayrı ayrı geçerlidir ve karışıklığa yol açmadığı yerlerde $i$ indisi yazılmamaktadır.
Eşitsizlikler gerçek sayı aritmetiğinde kesindir; gerçekleme bunları kW ya da kWh cinsinden, $\varepsilon_q$ payının kullanılmayan kısmının karşıladığı tek bir $10^{-9}$ toleransıyla denetler.

\subsection{Dikkat tabanlı politikanın boyuttan bağımsızlığı}\label{app:equivariance}

\begin{lemma}[Eşdeğişkenlik ve boyuttan bağımsızlık]\label{lem:equivariance}
$\varpi$, küresel belirteç sabit kalırken şarj ünitesi belirteçlerini ve maskelerini permüte etsin.
\begin{enumerate}
\item[(i)] Ortalamalar $\mu_{\theta}$ ve standart sapmalar $\sigma_{\theta}$ permütasyona eşdeğişkendir, değer $V_\omega$ ise permütasyona göre değişmezdir.
\item[(ii)] Maskelenmemiş belirteçlerin çıktıları ve değer, maskelenmiş belirteçlerin özniteliklerine ya da sayısına bağlı değildir.
\item[(iii)] Parametre sayısı $N$'ye bağlı değildir.
\end{enumerate}
\end{lemma}

\begin{proof}
Gömmeler, katman normalizasyonu, ileri beslemeli bloklar ve aktör başlığı satır bazında çalışır; dolayısıyla permütasyonla değişmeli (commute) davranır.
Tek bir dikkat başlığı için anahtar maskesi yalnızca anahtar indisine bağlıdır; bu nedenle belirteçleri ve maskeyi permüte etmek, küresel belirteci sabit bırakan $P$ permütasyon matrisi için maske matrisi $M$'yi $PMP^\top$'ye eşler.
Satır bazındaki softmax $\mathrm{softmax}(PXP^\top)=P\,\mathrm{softmax}(X)P^\top$ eşitliğini sağlar; dolayısıyla maskeli bir dikkat başlığı permüte edilmiş girdileri permüte edilmiş çıktılara eşler ve bileşke, aktör için (i)'yi verir.
Maskelenmiş bir anahtar her softmax'ta sıfır ağırlık alır; bu nedenle değeri hiçbir zaman başka bir belirtecin çıktısına girmez ve satır bazındaki yapıyla birlikte bu özellik (ii)'yi verir.
Dikkat havuzlama anahtarlar üzerinden toplam aldığından değişmezdir; her ağırlık matrisinin boyutu yalnızca $(F_c,F_g,d_m,n_H,n_L)$ ile belirlenir; dolayısıyla (iii) geçerlidir.
Güvenlik katmanı da permütasyona eşdeğişkendir, çünkü problemi $p$, $x$, $l$ ve $u$'nun ortak permütasyonu altında değişmezdir ve minimize edicisi tektir.
\end{proof}

Parametre sayıları aktör için 610.370, kritik için 676.481 olup her $N$ için toplam 1.286.851'dir.
Deep Sets, belirteç başına bir ağı her şarj ünitesi belirtecine, tüm belirteç gömmelerinin maskeli ortalaması ve küresel belirteçle birlikte uygular; dolayısıyla aynı argüman bu yapıyı da kapsar.
Genişliği, parametre sayısını \method{} ile \%10 içinde eşleştirecek biçimde seçilmiştir. 20 yuvalı ve 640 birimlik iki gizli katmanlı MLP karşılaştırma yönteminin aktörü 614.440 parametreye sahiptir; bu sayı yuva sayısıyla doğrusal olarak artar.

\subsection{Erişilebilir enerji ve eşikler}\label{app:reach}

\begin{lemma}[Erişilebilir enerji]\label{lem:reach}
A1 altında $\Psi_k$, her $k\ge0$ için sürekli ve azalmayandır.
$E(t)$'den başlayan herhangi bir kabul edilebilir güç dizisi $E(t+k)\le\Psi_k(E(t))$ sonucunu verir; eşitlik tam güçte şarj için sağlanır.
Bir EV, $t$ adımında bireysel olarak olurludur (yani kabul edilebilir bir kontrol \Eq{M-eq:requirement} koşulunu sağlar) ancak ve ancak $E(t)\ge B(\tau(t))$ ise.
\end{lemma}

\begin{proof}
Sürekli ve azalmayan eşlemelerin bileşkesi sürekli ve azalmayandır; dolayısıyla $\Psi_k=\varphi^{k}$ her iki özelliği de A1'den devralır.
Erişim sınırı için $k$ üzerinde tümevarım kullanılır; $k=0$ durumu açıktır.
Tümevarım adımında, kabul edilebilir bir $p\le\bar P(E)$ gücü, $g$ artan olduğundan $E(t+1)=E+\Delta t\,g(p)\le E+\Delta t\,g(\bar P(E))=\varphi(E)$ sonucunu verir; $E(t+1)$'den başlayan tümevarım hipotezi $\Psi_{k-1}$'in monotonluğuyla birlikte
\[
E(t+k)\le\Psi_{k-1}\big(E(t+1)\big)\le\Psi_{k-1}\big(\varphi(E)\big)=\Psi_k(E)
\]
sonucunu verir.
Tam güçte şarj her adımda eşitlik sağlar.

Eşik nitelendirmesi için tümevarımla $\{E\in[E^{\min},E^{\sm{cap}}]:\Psi_k(E)\ge\Ereq\}=[B(k),E^{\sm{cap}}]$ olduğu gösterilir; $k=0$ durumu $B(0)=\Ereq$ tanımıdır.
\Eq{M-eq:thresholds} içindeki küme boş değildir, çünkü $\varphi(E^{\sm{cap}})=E^{\sm{cap}}\ge\Ereq\ge B(k-1)$'dir.
Tümevarım adımında $\Psi_k(E)=\Psi_{k-1}(\varphi(E))\ge\Ereq$ ancak ve ancak $\varphi(E)\ge B(k-1)$ ise sağlanır; $\varphi$ sürekli ve azalmayan olduğundan $\{E:\varphi(E)\ge B(k-1)\}$ minimumu $B(k)$ olan kapalı bir üst kümedir.
Erişim sınırına göre, $E(t)$'den başlayarak kabul edilebilir bir kontrolün \Eq{M-eq:requirement} koşulunu sağlaması ancak ve ancak tam güçte şarjın bunu sağlamasıyla, yani ancak ve ancak $\Psi_{\tau(t)}(E(t))\ge\Ereq$ olmasıyla mümkündür.
\end{proof}

Sertifikalı bir tablo için \Eq{M-eq:certified} içindeki ilk iki özellik bileşen bazında $\Bh\ge B$ sonucunu doğurur, çünkü $\varphi(\Bh(k))\ge\Bh(k-1)\ge B(k-1)$ eşitsizliği \Eq{M-eq:thresholds} içindeki minimallik nedeniyle $\Bh(k)\ge B(k)$ sonucunu verir.

\subsection{İzdüşüm}\label{app:projection}

\begin{proof}[İspat (Önerme~\ref{M-prop:projection})]
(i) $p\in\mathcal K(t)$ ise $l\le p\le u$ ifadesi $l\le u$ ve $\sum_il_i\le\mathbf 1^\top p\le\sum_iu_i$ sonuçlarını verir; bu eşitsizlikler $-\Crev\le\mathbf 1^\top p\le\Ctil$ ile birleştirildiğinde dört koşulu verir.
Tersine, $\upsilon\in[0,1]$ için $p(\upsilon)=l+\upsilon(u-l)$ kutunun içinde kalır ve $\mathbf 1^\top p(\upsilon)$, $[\sum_il_i,\sum_iu_i]$ aralığındaki her değeri alır; bu aralık dört koşul altında $[-\Crev,\Ctil]$ ile kesişir.

(ii) \Eq{M-eq:proj} içindeki problem, boş olmayan bir çokyüzlü üzerinde kesin dışbükey bir kuadratiği en küçükler; dolayısıyla minimize edici vardır ve tektir, KKT koşulları ise gerekli ve yeterlidir.
Toplam kısıtları için $\mu^{+},\mu^{-}\ge0$ ve kutu için $\nu^{l},\nu^{u}\ge0$ çarpanlarıyla bu koşullar şöyledir:
\[
\begin{gathered}
p-x+(\mu^{+}-\mu^{-})\mathbf 1-\nu^{l}+\nu^{u}=0,\qquad \mu^{+}(\mathbf 1^\top p-\Ctil)=0,\qquad \mu^{-}(\mathbf 1^\top p+\Crev)=0,\\
\nu^{l}_i(p_i-l_i)=0,\qquad \nu^{u}_i(u_i-p_i)=0.
\end{gathered}
\]
Toplam yanıt $\mathcal S(\lambda)=\sum_i\clip(x_i-\lambda,l_i,u_i)$ sürekli, artmayan ve parçalı doğrusaldır; $\lambda\le\min_i(x_i-u_i)$ için $\mathcal S(\lambda)=\sum_iu_i$ ve $\lambda\ge\max_i(x_i-l_i)$ için $\mathcal S(\lambda)=\sum_il_i$'dir.
$\mathcal S(0)>\Ctil\ge\sum_il_i$ ise ara değer teoremi $\mathcal S(\lambda)=\Ctil$ denkleminin bir kökünü verir ve böyle her kök pozitiftir; $\mathcal S(0)<-\Crev\le\sum_iu_i$ durumu negatif bir kökle simetriktir.
Şimdi $p=\clip(x-\lambda^{*},l,u)$, $\mu^{+}=\max(\lambda^{*},0)$ ve $\mu^{-}=\max(-\lambda^{*},0)$ alınsın ve kutu çarpanları koordinat koordinat seçilsin.
$l_i<x_i-\lambda^{*}<u_i$ olan serbest koordinatlar $\nu_i^{l}=\nu_i^{u}=0$, $x_i-\lambda^{*}\le l_i$ olan koordinatlar $\nu_i^{l}=l_i-x_i+\lambda^{*}\ge0$ ve $\nu_i^{u}=0$, geri kalanlar ise $\nu_i^{u}=x_i-\lambda^{*}-u_i\ge0$ ve $\nu_i^{l}=0$ değerlerini alır.
Bu durumda durağanlık koordinat bazında sağlanır; tümleyenlik de sağlanır, çünkü $\lambda^{*}>0$ yalnızca $\mathbf 1^\top p=\Ctil$ olduğunda ve $\lambda^{*}<0$ yalnızca $\mathbf 1^\top p=-\Crev$ olduğunda gerçekleşir.
$\lambda^{*}=0$ için $\mathbf 1^\top p=\mathcal S(0)$ olan nokta olurludur; dolayısıyla her durumda $p=p^{*}$'dır.
$\lambda_1<\lambda_2$ her ikisi de kök ise her $\clip(x_i-\lambda,l_i,u_i)$ koordinatı $\lambda$'da artmayandır ve toplamlar eşittir; dolayısıyla tüm koordinatlar eşittir.
\end{proof}

PPO güvenlik katmanını ortamın bir parçası olarak ele alır ve Jacobi matrisini hiçbir zaman kullanmaz.
Bununla birlikte Jacobi matrisi, katmanın ham eylemdeki değişiklikleri komut edilen güçlere nasıl aktardığını betimler.

\begin{proposition}[İzdüşümün Jacobi matrisi]\label{prop:jacobian}
$\mathcal F=\{i:l_i<x_i-\lambda^{*}<u_i\}$ serbest küme ve $D$ bu kümenin göstergelerinden oluşan köşegen matris olsun.
$p^{*}=\Pi_t(x)$'in türevlenebilir olduğu her yerde Jacobi matrisi, toplam kısıtı etkin değilse $D$, $\mathcal F\neq\emptyset$ iken etkinse $D-\mathbf 1_{\mathcal F}\mathbf 1_{\mathcal F}^\top/|\mathcal F|$'dir; burada $\mathbf 1_{\mathcal F}$, $\mathcal F$'nin gösterge vektörüdür.
\end{proposition}

\begin{proof}
Toplam kısıtı etkin değilse yerel olarak $i\in\mathcal F$ için $p_i=x_i$'dir ve diğer durumlarda $p_i\in\{l_i,u_i\}$ sabittir; dolayısıyla Jacobi matrisi $D$'dir.
Toplam kısıtı $\vartheta\in\{\Ctil,-\Crev\}$ hedefinde etkin olduğunda ise $\sum_{i\in\mathcal F}(x_i-\lambda)+\sum_{i\notin\mathcal F}p_i=\vartheta$ olur.
Buradan $\partial\lambda/\partial x_j=\mathbb 1[j\in\mathcal F]/|\mathcal F|$ ve $\partial p_i/\partial x_j=\mathbb 1[i\in\mathcal F]\big(\mathbb 1[i=j]-\mathbb 1[j\in\mathcal F]/|\mathcal F|\big)$ elde edilir.
$\mathcal F=\emptyset$ ise $p$ noktası yerel olarak sabittir; $p^{*}$'ın türevlenebilir olmadığı sıfır ölçülü kümede ise komşu bölgelerin matrisleri Clarke'ın genelleştirilmiş Jacobi matrisinin elemanlarıdır.
\end{proof}

Jacobi matrisi simetrik ve idempotenttir; dışbükey bir küme üzerine Öklid izdüşümü olarak $\Pi_t$ genişletmeyen (non-expansive) bir eşlemedir.
Etkin bir satırda Jacobi matrisi tüm kırpılmış koordinatlar için sıfırdır; bu, katmanı ortamın parçası olarak ele almayı gerekçelendiren sıfır gradyan etkisidir \cite{kasaura2023benchmarking}.

\subsection{Tek adımlık değişmezlik}\label{app:onestep}

Ana metindeki Teorem~\ref{M-thm:invariance} için verilen ispat taslağı ve aşağıdaki tam ispat, her ikisi de aşağıdaki lemmaya dayanmaktadır.

\begin{lemma}[Tek adımlık değişmezlik]\label{lem:onestep}
A1 ve A3 varsayılsın ve $\Bh$, \Eq{M-eq:certified} koşullarını sağlasın.
EV $i$, $t$ adımında $\tau=\tau_i(t)$ kalan eylem ve $E=E_i(t)\ge\Bh(\tau)$ enerjisiyle bağlı olsun.
Bu durumda $\pmn(E)\le l_i(t)\le u_i(t)$'dir ve komut edilen her $p\in[l_i(t),u_i(t)]$ gücü $E_i(t+1)\ge\Bh(\tau-1)$ sonucunu doğurur.
\end{lemma}

\begin{proof}
$\pmn\le l_i$ sınırı geçerlidir, çünkü $\pmn\le0\le u_i$'dir ve \Eq{M-eq:bounds} içindeki dış minimumun her iki argümanı da en az $\pmn$'dir.
$l_i\le u_i$ için dış minimum yeterlidir ve koruma kesin aritmetikte etkin değildir; çünkü $g^{-1}\big((\Bh(\tau-1)-E)/\Delta t\big)\le\bar P(E)$, $\Bh(\tau-1)\le\varphi(E)$ ile denktir ve A1, \Eq{M-eq:certified} ile birlikte $\varphi(E)\ge\varphi(\Bh(\tau))\ge\Bh(\tau-1)$ sonucunu verir.
İkinci iddiaya geçersek, $y=(\Bh(\tau-1)-E)/\Delta t$ olsun; kesin aritmetikte $l_i\ge\min\{u_i,g^{-1}(y)\}=g^{-1}(y)$ olduğundan ve $g$ artan olduğundan her $p\ge l_i$ için $g(p)\ge y$ sağlanır ve A3, $E_i(t+1)\ge E+\Delta t\,g(p)\ge\Bh(\tau-1)$ sonucunu verir.
\end{proof}

$\tau\ge2$ olduğunda EV $t+1$ adımında bireysel olarak olurludur, çünkü $\tau_i(t+1)=\tau-1$'dir; $\tau=1$ olduğunda ise lemma $E_i(d_i+1)\ge\Bh(0)\ge\Ereq_i$ sonucunu verir.

\subsection{Olurluluğun korunması}\label{app:theorem}

Teorem~\ref{M-thm:invariance} için verilen ispat sırasıyla olurlu adımları, olursuz adımları ve ayrılışa ilişkin kısmı ele almaktadır.
(a) iddiası yalnızca A3 ve A6'yı kullanır; ayrılışa ilişkin kısım ise ek olarak A1, A2 ve A4'ü kullanır.
A1--A6 koşulları \Sec{M-sec:invariance} içinde listelenenlerdir; A5 olurlu adımın tanımıdır.

\begin{proof}[İspat (Teorem~\ref{M-thm:invariance})]
\emph{Olurlu adımlarda (a) iddiası.}
\Eq{M-eq:bounds} içindeki koruma, EV bireysel olarak olurlu olsun ya da olmasın, her adımda bağlı her EV için $\pmn_i\le l_i\le u_i$ sonucunu verir; boş şarj ünitelerinde ise $l_i=u_i=0$'dır.
\Alg{alg:safety} içindeki olurluluk testi geçilirse Önerme~\ref{M-prop:projection}(i) $\mathcal K(t)\neq\emptyset$ sonucunu verir; dolayısıyla her sonlu $x(t)$ için $p(t)=\Pi_t(x(t))\in\mathcal K(t)$'dir.
Fiziksel limitler $p_i(t)\in[l_i,u_i]\subseteq[\pmn_i(E_i(t)),\bar P_i(E_i(t))]$ ilişkisinden çıkar.
Trafo limiti için $\mathbf 1^\top p(t)\in[-\Crev(t),\Ctil(t)]$'dir ve A3
\[
\sum_i\tilde p_i(t)\in\Big[\sum_ip_i(t),\ \sum_ip_i(t)+n(t)\varepsilon_q\Big]\subseteq\big[-\Crev(t),\ \Ctil(t)+n(t)\varepsilon_q\big]=\big[-\Crev(t),C(t)\big]
\]
sonucunu verir.

\emph{Olursuz adımlarda (a) iddiası.}
A6 altında $\Ctil(t)=C(t)-n(t)\varepsilon_q\ge(N-n(t))\varepsilon_q\ge0$ ve $\Crev(t)\ge0$'dır; dolayısıyla yedek kuralın (i) durumu $\Ctil<-\Crev\le0$ gerektirir ve bu nedenle olanaksızdır.
$u_i\ge0$ olduğundan $\sum_iu_i<-\Crev\le0$ olan (ii) durumu da olanaksızdır.
(iv) alt durumu $b=\Ctil-\sum_i\pmn_i<0$ gerektirir ve $\pmn_i\le0$ olduğundan ortaya çıkamaz.
Dolayısıyla her olursuz adım (iii) durumuna ait bir adımdır ve Lemma~\ref{lem:fallback} $\pmn\le p\le l\le u$ ile $\mathbf 1^\top p=\Ctil\in[-\Crev,\Ctil]$ sonuçlarını verir.
Fiziksel limitler buradan çıkar ve A3, toplam sınırını olurlu adımlarda olduğu gibi gerçekleşen güçlere aktarır.

\emph{(b) iddiası.}
Kalış süresinin her adımında $p_i(t)\ge l_i(t)$ olan bir EV $i$ sabitlensin ve $E_i(t)\ge\Bh_i(\tau_i(t))$ değişmezi ele alınsın.
Varışta bu değişmez A4'tür.
Değişmez $t$ adımında sağlanıyorsa $l_i(t)$, \Eq{M-eq:bounds} içindeki birinci durumu izler ve izdüşümün ya da yedek kuralın her çıktısı $p_i(t)\le u_i(t)$'yi sağlar; dolayısıyla $p_i(t)\in[l_i(t),u_i(t)]$'dir.
Bu durumda Lemma~\ref{lem:onestep} $E_i(t+1)\ge\Bh_i(\tau_i(t)-1)$ sonucunu verir.
$d_i>t$ için bu eşitsizlik $t+1$ adımındaki değişmezdir; $d_i=t$ için ise $E_i(d_i+1)\ge\Bh_i(0)\ge\Ereq_i$ ayrılış gereksinimidir.
Argüman yalnızca EV $i$'ye gönderilen komutları kullanır; dolayısıyla diğer şarj ünitelerinde ne olursa olsun geçerlidir.
Kalış süresi boyunca olursuz adımla karşılaşmayan her EV, kalışının her adımında $p_i\in[l_i,u_i]$ alır ve kapsama dâhildir.
EV2Gym her EV'yi en geç $T-2$ adımında çıkardığından epizottaki her EV ayrılışına ufuk içinde ulaşır ve argümanın ayrılışa ilişkin kısmı bunların tümüne uygulanır.

\emph{(c) iddiası.}
A5 her adımda sağlanıyorsa epizodun hiçbir adımı olursuz değildir ve dolayısıyla (b) iddiası epizottaki her EV'ye uygulanır.
\end{proof}

A2, sınırların her adımda ihtiyaç duyduğu kalan eylemler $\tau_i(t)$ aracılığıyla devreye girer.
Kabul kontrolü varıştan sonra bir gereksinimi hiçbir zaman değiştirmez; dolayısıyla değişmezde kullanılan gereksinim kalış süresince sabittir.

\textbf{A5 için politikadan bağımsız bir yeterli koşul.}
Alt sınır $l_i$, birinci durumundaki her iki terim de $E$'de artmayan olduğundan $E$'de artmayandır; $l_i^{\max}(t)$ ile $l_i$'nin $E=\Bh_i(\tau_i(t))$ noktasındaki değeri gösterilsin.
Her adımda $\Crev(t)\ge0$, $\Ctil(t)\ge-\Crev(t)$ ve $\sum_{i\in\mathcal N(t)}l_i^{\max}(t)\le\Ctil(t)$ olduğu varsayılsın.
Değişmez altında $l_i(t)\le l_i^{\max}(t)$ ve $\sum_iu_i\ge0\ge-\Crev$'dir; dolayısıyla A5 $t$ adımında sağlanır ve tümevarımla A5 her politika için sağlanır.
Ancak $\Bh_i(\tau_i)>E_i^{\min}$ olduğunda $l_i^{\max}(t)=\bar P_i(\Bh_i(\tau_i))$'dir; dolayısıyla bu koşul, trafonun önemsiz olmayan biçimde kısıtlanmış tüm EV'lere aynı anda tam güçte hizmet vermesini gerektirir.

\subsection{Yedek kural ve açık sınırı}\label{app:fallback}

Olurluluk testi başarısız olduğunda izdüşümün yerini yedek kural alır.
$l\le u$ yapısı gereği sağlandığından üç durum kalır; bunlardan üçüncüsünün bir alt durumu vardır.
\begin{enumerate}
\item[(i)] Çelişkili limitler, $\Ctil<-\Crev$: (iii) durumunun prosedürü uygulanır; dolayısıyla ileri yön limiti önceliklidir.
\item[(ii)] Aşırı PV ihracatı, $\sum_iu_i<-\Crev$: dolu her şarj ünitesi $p_i=u_i$ alır.
\item[(iii)] Gerekli güç fazlası, $\sum_il_i>\Ctil$: her EV $\pmn_i$'den başlar ve $b=\Ctil-\sum_i\pmn_i$ bütçesi esneklik payı artan sırada dağıtılır; her EV en fazla $l_i-\pmn_i$ kadar alır.
\item[(iv)] (i) ve (iii) durumlarının $b<0$ olan alt durumu: çıktı $p=\pmn$'de, yani ileri yön limitine en yakın kabul edilebilir komutta kalır.
\end{enumerate}
(iii) durumunda eşitlik hâlinde önce daha büyük $l_i-\pmn_i$ farkı, ardından daha küçük şarj ünitesi indisi öncelik alır.
Kayıtlar (i), (ii), (iii) ve (iv) durumlarını 3, 2, 1 ve 0 sayısal kodlarıyla FB-1*, FB-2, FB-1 ve FB-0 olarak etiketler ve her analiz tablosu bu eşlemeyi kullanır.
A6 altında (i), (ii) ve (iv) durumları ortaya çıkamaz (Teorem~\ref{M-thm:invariance} ispatı); dolayısıyla sonuç tablolarındaki sayıları her tekrarda A6'yı doğrular.

\begin{lemma}[Yedek kural]\label{lem:fallback}
Yedek kuralın $b\ge0$ olan (iii) durumunda çıktı $\pmn\le p\le l$ ve $\mathbf 1^\top p=\Ctil$ koşullarını sağlar; dolayısıyla trafo limiti ve fiziksel limitler sağlanır.
$t$ adımında bireysel olarak olurlu olan ve $p_i=l_i$ alan her EV bireysel olarak olurlu kalır.
(ii) ve (iv) durumlarında çıktı, $[\pmn,u]$ kutusu üzerinde trafo ihlalini en küçükler; dolayısıyla hiçbir kabul edilebilir komut bu ihlalden kaçınamaz.
\end{lemma}

\begin{proof}
(iii) durumunda $\sum_i(l_i-\pmn_i)>\Ctil-\sum_i\pmn_i=b$'dir; dolayısıyla bütçe, her EV $l_i$'ye ulaşmadan tükenir; buradan $\mathbf 1^\top p=\Ctil$ ve $p\le l$ elde edilir.
(i) durumu dışında ayrıca $\Ctil\ge-\Crev$'dir ve $p_i=l_i$ alan bireysel olarak olurlu bir EV, Lemma~\ref{lem:onestep} gereğince olurlu kalır.
$[\pmn,u]$ kutusu üzerinde $\mathbf 1^\top p$ toplamı $\pmn$'de en küçük, $u$'da en büyüktür; dolayısıyla (iv) ve (ii) durumlarının çıktıları komut edilen ihlali en küçükler.
\end{proof}

\begin{lemma}[Açık sınırı]\label{lem:deficit}
A1, A3 ve A4 varsayılsın ve $\bar P_i$'nin $E$'de artmayan olduğu varsayılsın.
EV2Gym gerçeklememizdeki gibi $\eta^{\sm c}=\eta^{\sm d}=1$ ile filonun bir epizottaki toplam ayrılış eksikliği en fazla
\begin{equation}\label{eq:deficit}
\Delta t\sum_{t\in\mathcal T_{\mathrm{iii}}}\Big(\sum_il_i(t)-\Ctil(t)\Big)+\Delta t\sum_{t\in\mathcal T_{\mathrm{iv}}}\sum_i\big(l_i(t)-\pmn_i(t)\big)
\end{equation}
kadardır; burada $\mathcal T_{\mathrm{iii}}$ ve $\mathcal T_{\mathrm{iv}}$, (iii) ve (iv) durumlarına ait adımlardır.
\end{lemma}

\Sec{sec:ev2gym} içindeki Lemma~\ref{lem:taper}, EV2Gym'in şarj gücü azalmasının Lemma~\ref{lem:deficit} içindeki ek koşulu sağladığını göstermektedir.

\begin{proof}
$E_i(t)\ge\Bh_i(\tau_i(t))$ ise $\tilde l_i(t)=l_i(t)$, aksi hâlde $\tilde l_i(t)=u_i(t)$ olsun ve $\tau(d_i+1)=0$ olmak üzere olurluluk açığı $\Delta E^{\sm{def}}(t)=(\Bh(\tau(t))-E(t))^{+}$ olarak tanımlansın.
$\Bh(0)\ge\Ereq$ olduğundan eksik enerji $(\Ereq-E(d_i+1))^{+}$ en fazla $\Delta E^{\sm{def}}(d_i+1)$'dir ve A4 gereğince $\Delta E^{\sm{def}}(t_i^{\sm{arr}})=0$'dır.
Dolayısıyla her adımda $\Delta E^{\sm{def}}(t+1)\le\Delta E^{\sm{def}}(t)+\Delta t\,\big(g(\tilde l)-g(p)\big)^{+}$ olduğunu göstermek ve bu artışları kalış süresi boyunca toplamak yeterlidir.
$E\ge\Bh(\tau)$ ise \Eq{M-eq:bounds} gereğince $E+\Delta t\,g(l)\ge\Bh(\tau-1)$'dir ve A3, $E(t+1)\ge\Bh(\tau-1)-\Delta t\,(g(l)-g(p))^{+}$ sonucunu verir.
$E<\Bh(\tau)$ olduğunda ise $\tilde l=u$ ve $E+\Delta t\,g(u)=\varphi(E)$'dir; \Eq{M-eq:certified}, $\varphi$'nin 1-Lipschitz özelliğiyle birlikte $\Bh(\tau-1)-\varphi(E)\le\varphi(\Bh(\tau))-\varphi(E)\le\Bh(\tau)-E=\Delta E^{\sm{def}}(t)$ sonucunu verir.
Dolayısıyla bu durumda da $\Bh(\tau-1)-E(t+1)\le\Delta E^{\sm{def}}(t)+\Delta t\,(g(u)-g(p))$'dir.
$g$ özdeşlik fonksiyonu olmak üzere filo sınırı için olurlu adımlarda $p_i\ge\tilde l_i$, yedek kuralın (ii) durumunda ise $p=u$ sağlanır.
(iii) durumunda $p\le l$ ve $\sum_i(l_i-p_i)=\sum_il_i-\Ctil$'dir; (iv) durumunda $p=\pmn$'dir; dolayısıyla adım başına artışların tüm EV'ler üzerinden toplanması \Eq{eq:deficit} ifadesini verir.
\end{proof}

\subsection{Konformal ayrılış garantisi}\label{app:conformal}

\begin{proof}[İspat (Önerme~\ref{M-prop:conformal})]
(i) Değiştirilebilirlik gereğince test skoru $e_{n_{\sm{cal}}+1}=\hat q(\mathbf w_i)-d_i$'nin $e_1,\dots,e_{n_{\sm{cal}}+1}$ skorları arasındaki sırası, eşitlikler dışında düzgün dağılımlıdır; dolayısıyla $\mathbb P(e_{n_{\sm{cal}}+1}\le\hat Q_{1-\alpha})\ge\lceil(n_{\sm{cal}}+1)(1-\alpha)\rceil/(n_{\sm{cal}}+1)\ge1-\alpha$'dır \cite{vovk2022algorithmic,lei2018distributionfree}.
$e_{n_{\sm{cal}}+1}\le\hat Q_{1-\alpha}$ olayı $d_i\ge\hat q(\mathbf w_i)-\hat Q_{1-\alpha}$ ile, $d_i$ bir tam sayı olduğundan da $d_i\ge\lceil\hat q(\mathbf w_i)-\hat Q_{1-\alpha}\rceil$ ile denktir.
$t_i^{\sm{arr}}$ ile maksimumun alınması hiçbir şeyi değiştirmez, çünkü $d_i\ge t_i^{\sm{arr}}$ kesinlikle sağlanır.

(ii) $d_i\ge\hat d_i^{\sm{lo}}$ olduğu ve kalış süresinin her adımında $p_i(t)\ge l_i(t)$ olduğu olay ele alınsın.
Bu olay üzerinde, $t=\hat d_i^{\sm{lo}}$'ya kadar $\tau^{\sm{lo}}$ ile uygulanan Lemma~\ref{lem:onestep}, $E_i(\hat d_i^{\sm{lo}}+1)\ge\Bh(0)$ sonucunu verir.
Kalan $\hat d_i^{\sm{lo}}<t\le d_i$ adımları için $\tau^{\sm{lo}}=1$'dir ve \Eq{M-eq:certified} içindeki üçüncü özellik $\Bh(1)\le\max\{\Bh(0),E^{\min}\}\le E_i(t)$ sonucunu verir.
Bu durumda $\tau=1$ ile Lemma~\ref{lem:onestep} yine $E_i(t+1)\ge\Bh(0)\ge\Ereq$ sonucunu verir; dolayısıyla olay üzerinde $E_i(d_i+1)\ge\Ereq$ sağlanır ve birleşim sınırı ifade edilen olasılığı verir.
Kalış süresinin hiçbir adımı hemen hemen kesin olarak olursuz değilse ikinci olasılık sıfırlanır ve sınır $1-\alpha$ olur.
Sınır, A5'in sağlandığı olaya koşullu olarak geçerli değildir, çünkü bu olay $d_i$ ile ilişkili olabilir.
\end{proof}

% ---------------------------------------------------------------------
\section{EV2Gym modellemesine ilişkin notlar}\label{sec:ev2gym}
% ---------------------------------------------------------------------

Bu bölüm, ana metindeki modelin \texttt{55ef368c} commit'indeki EV2Gym'e nasıl karşılık geldiğini belgelemektedir; simülatör kaynak kodunda değişiklik yapılmadan kullanılmaktadır.
Simülatörün varsayılan ayarlarından tüm sapmalar kendi sarmalayıcımız, yapılandırma dosyalarımız ve EV belirtim dosyamız aracılığıyla yapılmaktadır.

\subsection{Şarj gücü azalması ve A1 varsayımı}\label{app:taper}

Şarj durumu birimlerinde $\varsigma=E/E^{\sm{cap}}$ olmak üzere $\delta=\Pnom\Delta t/E^{\sm{cap}}$ tam güçteki bir adımın kazancı, $\varsigma^{\sm{tr}}\in(0,1]$ geçiş şarj durumu, $m$ şarj gücü azalması çarpanı ve $\kappa=m\delta/(1-\varsigma^{\sm{tr}})$ olsun.
EV'nin en yüksek alternatif akım şarj gücü şarj ünitesinin anma değerini aşmadığında, EV2Gym'in şarj rutininin tam akımda okunması şu sonucu verir:
\begin{equation}\label{eq:taper}
\frac{\varphi(E)}{E^{\sm{cap}}}=\min\{\hat\varphi(\varsigma),\varsigma+\delta\},\qquad
\hat\varphi(\varsigma)=\begin{cases}
\varsigma+\delta, & \varsigma+\delta\le\varsigma^{\sm{tr}},\\
1-(1-\varsigma^{\sm{tr}})\exp\!\Big(-m\dfrac{\varsigma+\delta-\varsigma^{\sm{tr}}}{1-\varsigma^{\sm{tr}}}\Big), & \varsigma<\varsigma^{\sm{tr}}<\varsigma+\delta,\\
1-(1-\varsigma)\,e^{-\kappa}, & \varsigma\ge\varsigma^{\sm{tr}}.
\end{cases}
\end{equation}
Kodla birlikte yayımlanan bir denetim betiğine göre bu kapalı biçim, EV2Gym'in şarj rutiniyle $1.4\times10^{-14}$ kWh duyarlılığında örtüşmektedir.
Denetim; EV modeli, verim, geçiş şarj durumu, $m\in\{1,5\}$ çarpanı ve $\Delta t\in\{5,15\}$ dk adım uzunluğunun 768 kombinasyonunu kapsamaktadır.
\TODO{Yazarlar bu denetimi yeniden çalıştırıp doğrular}.

\begin{lemma}[EV2Gym'de A1]\label{lem:taper}
EV'nin en yüksek alternatif akım şarj gücünün şarj ünitesinin anma değerini aşmadığını varsayalım.
Bu durumda her adım uzunluğu için EV2Gym'in \Eq{eq:taper} ile verilen tek adımlık tam güç eşlemesi sürekli, azalmayan ve 1-Lipschitz'tir.
\end{lemma}

\begin{proof}
$\varsigma=\varsigma^{\sm{tr}}-\delta$ noktasında $\hat\varphi$'nin ilk iki dalı da $\varsigma^{\sm{tr}}$'ye, $\varsigma=\varsigma^{\sm{tr}}$ noktasında ise son iki dalı da $1-(1-\varsigma^{\sm{tr}})e^{-\kappa}$'ya eşittir.
Dolayısıyla eşleme, sürekli fonksiyonların minimumu olarak süreklidir.
Dallarının eğimleri $1$, $m\exp(\cdot)>0$ ve $e^{-\kappa}\in(0,1)$'dir ve $\varsigma+\delta$'nın eğimi 1'dir; dolayısıyla eşleme, azalmayan fonksiyonların minimumu olarak azalmayandır ve A1 sağlanır.
Yalnızca orta dalın eğimi 1'in üzerinde olabilir; bu da $m>1$ olduğunda gerçekleşir.
$(\varsigma^{\sm{tr}}-\delta,\varsigma^{\sm{tr}})$ aralığında $\hat h(\varsigma)=\hat\varphi(\varsigma)-\varsigma-\delta$ fonksiyonu kesin içbükeydir ve $\hat h(\varsigma^{\sm{tr}}-\delta)=0$'dır.
Bir iç noktada $\hat h(\varsigma_1)<0$ ise içbükeylik $\hat h'(\varsigma_1)\le \hat h(\varsigma_1)/(\varsigma_1-\varsigma^{\sm{tr}}+\delta)<0$ sonucunu verir; dolayısıyla $\hat\varphi'(\varsigma_1)<1$'dir.
Böylece \Eq{eq:taper} içindeki minimumun kesin olarak $\hat\varphi$'yi seçtiği her yerde eğim 1'in altındadır; diğer yerlerde minimum, eğimi 1 olan $\varsigma+\delta$'yı seçer.
Dolayısıyla eşleme, eğimleri $[0,1]$ aralığında olan parçalı sürekli türevlenebilir bir fonksiyondur ve bu nedenle 1-Lipschitz'tir.
\end{proof}

EV2Gym'de $\eta^{\sm c}=1$ ve $\varphi(E)=E+\Delta t\,\bar P(E)$'dir; dolayısıyla 1-Lipschitz bir $\varphi$, Lemma~\ref{lem:deficit} içindeki ek koşul olan artmayan bir $\bar P$ ile denktir.
Gücü kWh başına $1/(\eta^{\sm c}\Delta t)$ kW'tan daha hızlı düşen bir şarj gücü azalması A1'i ihlal eder ve sertifikalı eşikler ayrılış sonucunu onaramaz.

\subsection{Eyleyici davranışı ve A3 varsayımı}\label{app:actuation}

Karar değişkeni $p_i$, EV2Gym'in raporladığı, faturalandırdığı ve trafoya gönderdiği güçtür ve adım başına batarya tarafındaki enerji değişimine eşittir.
Bu kabulle \Eq{M-eq:dynamics} içinde $\eta_i^{\sm c}=\eta_i^{\sm d}=1$ olur; EV2Gym'in şarj ve deşarj verimleri $\zeta_i^{\sm c}$ ve $\zeta_i^{\sm d}$ ise yalnızca güç limitleri üzerinde düşürme (derating) çarpanları olarak etki eder.

Belirtim dosyamız her EV modeli için tek bir verim sabitlemektedir: yedi model için 0.90, bir model için 0.93; bu değerler simülatörün tablosundan modelin nominal akımında alınmıştır.
EV2Gym, normalize edilmiş eylemi (komutun $\zeta_i^{\sm c}P^{\sm{cs}}$'ye bölümü; burada şarj ünitesi anma değeri $P^{\sm{cs}}=22.17$ kW'tır) beş ondalık basamağa yuvarlar ve 0.01 A'in altındaki deşarj akımlarını yok sayar.
Ayrıca her etkin adımdan sonra depolanan enerjiyi, eklenen enerjiyi faturalandırmadan 0.01 kWh'e yukarı yuvarlar.
Bu nedenle eyleyici eşlemesi, bir komutu $u_i$'nin $\varepsilon_{\sm{sat}}=10^{-3}$ kW yakınında olduğunda tam akım olarak gönderir; aksi hâlde komutu beş ondalık basamaklı eylem ızgarasında yukarı yuvarlar.
\Alg{alg:safety}, bu eyleyici davranışıyla güvenlik katmanının tam bir adımını listelemektedir.
$\sum_il_i$ makine duyarlılığında $\Ctil$ değerine eşit olduğunda yuvarlama toplam limiti bozabilir.
Bu nedenle gerçekleme her izdüşümden sonra limiti denetler ve başarısız olan bir satırı, aralığının olurlu ucunu koruyan bir ikiye bölme yöntemiyle yeniden hesaplar.

\begin{algorithm}[t]
\caption{Güvenlik katmanının bir adımı}\label{alg:safety}
\begin{algorithmic}[1]
\Require ham eylem $x\in\R^N$ (sonlu olmayan girdiler 0 yapılır); bağlı her EV için: $E_i(t)$, $\tau_i(t)$, sertifikalı tablo $\Bh_i$; limitler $C(t)$, $\Crev(t)$
\For{dolu her şarj ünitesi $i$}
  \State $u_i\gets\bar P_i(E_i)$;\quad $\pmn_i\gets$ \Eq{M-eq:limits}
  \If{$E_i\ge\Bh_i(\tau_i)-10^{-9}$} $l_i\gets\min\{u_i,\max\{\pmn_i,\,g_i^{-1}((\Bh_i(\tau_i-1)-E_i)/\Delta t)\}\}$ \Else\ $l_i\gets u_i$ \EndIf
  \State $\chi_i\gets(E_i-\Bh_i(\tau_i))/(\eta_i^{\sm c}\Pnom_i\Delta t)$
\EndFor
\State $\Ctil\gets C-n(t)\varepsilon_q$
\If{$\sum_il_i\le\Ctil$ \textbf{and} $\sum_iu_i\ge-\Crev$ \textbf{and} $\Ctil\ge-\Crev$} \Comment{olurlu adım (A5)}
  \State $p^{0}\gets\clip(x,l,u)$
  \If{$\mathbf 1^\top p^{0}\in[-\Crev,\Ctil]$} $p\gets p^{0}$
  \Else
    \State $\{x_i-u_i\}\cup\{x_i-l_i\}$ kırılma noktalarını sırala; $\lambda^{*}$ değerini tek geçişte bul
    \State $p\gets\clip(x-\lambda^{*},l,u)$; $\mathbf 1^\top p\notin[-\Crev,\Ctil]$ ise güvenli ikiye bölme ile yeniden hesapla
  \EndIf
\Else
  \State $p\gets$ \Sec{app:fallback} içindeki yedek kural; durumu kaydet
\EndIf
\State $p$'yi simülatöre gönder: $p_i\ge u_i-10^{-3}$ ise tam akım, değilse eylem ızgarasında yukarı yuvarlanmış değer
\State \Return $p$
\end{algorithmic}
\end{algorithm}

\begin{lemma}[EV2Gym'de A3]\label{lem:a3}
Her EV'nin en yüksek alternatif akım şarj gücünün şarj ünitesinin anma değerini aşmadığını ve verimlerin akım düzeylerinden bağımsız olarak sabit olduğunu varsayalım.
Şarj üniteleri ve EV'ler üç fazlı olsun; her şarj ünitesinde bir EV bulunsun ve şarj için ölü bantlar sıfır olsun.
Ayrıca V2G'nin etkin olduğunu ve tüm EV'lerin epizot bitmeden ayrıldığını varsayalım; $m>1$ için $z_0(m)\in(0,1)$, $1-e^{-mz}=z$ denkleminin pozitif kökü olsun.
Her EV için $\kappa\le\kappa^{*}(m)=m\,z_0(m)$ ise eyleyici eşlemesi A3'ü $\varepsilon_q=0.008$ kW ile gerçekler.
\end{lemma}

\begin{proof}
Önce, normalize edilmiş eylemin pilot değeri $p/(\zeta^{\sm c}P^{\sm{cs}})$ olarak gönderilen en yüksek değerin altındaki her $p\in(0,\bar P(E)]$ şarj komutunun tam olarak teslim edildiği gösterilir.
$\hat\delta=p\Delta t/E^{\sm{cap}}$ bu komutun kazancı olsun; $\hat\delta\le\hat\delta^{*}=\varphi(E)/E^{\sm{cap}}-\varsigma\le\delta$'dır.
Düşürülmüş bir pilot için EV2Gym geçiş noktasını $\varsigma_x=1-(\hat\delta/\delta)(1-\varsigma^{\sm{tr}})\ge\varsigma^{\sm{tr}}$ değerine kaydırır ve $\Omega=1-\varsigma_x$ yazılır.
Teslim edilen kazanç $\min\{\hat\delta,\hat\varphi_x(\varsigma)-\varsigma\}$'dir; burada $\hat\varphi_x$, $(\varsigma^{\sm{tr}},\delta)$ yerine $(\varsigma_x,\hat\delta)$ konmuş \Eq{eq:taper} eşlemesidir ve kaydırılmış şarj gücü azalmasının üssü $m\hat\delta/\Omega$, $\kappa$'ya eşittir.

$\varsigma+\hat\delta\le\varsigma_x$ ise adım sabit akım dalında kalır ve kazanç $\hat\delta$'dır.
$\varsigma\ge\varsigma_x$ için azalma dalı $\hat\varphi_x(\varsigma)-\varsigma=(1-\varsigma)(1-e^{-\kappa})$ sonucunu verir.
$\varsigma\ge\varsigma^{\sm{tr}}$ olduğundan \Eq{eq:taper} $\hat\delta\le\hat\delta^{*}\le(1-\varsigma)(1-e^{-\kappa})$ sonucunu verir ve kazanç yine $\hat\delta$'dır.

Kalan $\varsigma<\varsigma_x<\varsigma+\hat\delta$ durumunda $y=\varsigma+\hat\delta-\varsigma_x>0$ ve $\psi=y/\Omega$ olsun.
Bu durumda $\hat\varphi_x(\varsigma)-\varsigma-\hat\delta=\Omega(1-e^{-m\psi}-\psi)$ olur; bu, $\psi$'nin 0'da ve $z_0(m)$'de sıfırlanan içbükey bir fonksiyonudur; dolayısıyla $\psi\le z_0(m)$ olduğunda kazanç $\hat\delta$'dır.
$\psi=1-(1-\varsigma-\hat\delta)/\Omega$ olduğundan bu koşul $1-\varsigma-\hat\delta\ge(1-z_0)\,\hat\delta(1-\varsigma^{\sm{tr}})/\delta$ biçimini alır.
$\varsigma<\varsigma^{\sm{tr}}$ ise \Eq{eq:taper} gereğince $1-\varsigma-\hat\delta\ge1-\varphi(E)/E^{\sm{cap}}\ge(1-\varsigma^{\sm{tr}})e^{-\kappa}$'dır ve $\hat\delta\le\delta$ ile koşul, $e^{-\kappa}\ge1-z_0=e^{-mz_0}$ ise, yani $\kappa\le mz_0$ ise sağlanır.
$\varsigma\ge\varsigma^{\sm{tr}}$ için $\hat\delta\le(1-\varsigma)(1-e^{-\kappa})$ sınırı $1-\varsigma-\hat\delta\ge(1-\varsigma)e^{-\kappa}$ sonucunu verir.
Bu durumda koşulun sağ tarafı en fazla $(1-z_0)(1-\varsigma)(1-e^{-\kappa})m/\kappa$'dır; dolayısıyla $\kappa/(e^{\kappa}-1)\ge m(1-z_0)$ olması yeterlidir.
$\kappa=mz_0$ noktasında iki taraf eşittir, çünkü $e^{mz_0}=1/(1-z_0)$'dır ve sol taraf $\kappa$'da azalır; dolayısıyla yine $\kappa\le mz_0$ yeterlidir.
Deşarjda şarj gücü azalması yoktur; EV2Gym, $[\pmn(E),0)$ içindeki bir deşarj komutunu $E^{\min}$ ile alttan ve deşarj limitiyle üstten sınırlanmış olarak tam teslim eder.

Ardından gerçekleşen güç sınırlanır.
Doymuş bir komut $\varphi(E)-E=\Delta t\,u_i$ teslim eder ve $u_i-p_i\le\varepsilon_{\sm{sat}}=10^{-3}$ kW'tır.
$u_i-\varepsilon_{\sm{sat}}$'nin altındaki her pozitif komut en fazla $\zeta^{\sm c}P^{\sm{cs}}10^{-5}\le2.22\times10^{-4}$ kW yukarı yuvarlanır, $u_i$'nin altında kalır ve ilk kısım gereğince tam teslim edilir; EV2Gym'in kendi yuvarlaması ise ızgara üzerindeki bir değeri değiştirmez.
Negatif komutlar en fazla $2.22\times10^{-4}$ kW sıfıra doğru yuvarlanır; ölü bant devreye girerse gerçekleşen güç 0 olur ve $0-p_i<0.0069+2.22\times10^{-4}$ kW'tır.
Aksi hâlde $\tilde p_i\ge p_i\ge\pmn(E)$'dir; dolayısıyla ne deşarj limiti ne de enerji tabanı komutu kırpar.
Her durumda $0\le\tilde p_i-p_i\le0.0072<\varepsilon_q$ kW'tır ve boş şarj üniteleri sıfır güç çeker.

Son olarak depolanan enerji $E(t+1)=\lceil E+\Delta t\,\tilde p_i\rceil_{0.01}\ge E+\Delta t\,\tilde p_i\ge E+\Delta t\,p_i$'dir.
Ayrıca $E+\Delta t\,\tilde p_i\le\varphi(E)\le E^{\sm{cap}}$'dir ve her kapasite, yuvarlamadan sonra tam olarak temsil edilebilen 0.01 kWh'in bir katı olduğundan yukarı yuvarlama $E^{\sm{cap}}$ değerini aşamaz.
\end{proof}

$m=5$ için kök $z_0(5)=0.99302$ ve eşik $\kappa^{*}(5)=4.965$'tir.
Yapılandırmamızda sekiz EV modelinin tümü için $\Delta t=15$ dk'da $\kappa\le2.391$ ve $\Delta t=5$ dk'da $\kappa\le0.797$'dir; dolayısıyla Lemma~\ref{lem:a3} uygulanabilir.
$m\le1$ için teslim, geçiş noktası yakınında kesin değildir; bu nedenle şarj gücü azalması çarpanı varsayılan değeri olan 5'te tutulmaktadır.
Kodla birlikte yayımlanan kapalı döngü bir denetime göre, ileri yön limiti $n(t)\varepsilon_q$ kadar sıkılaştırılmasaydı gerçekleşen güçler sıkı adımlarda ileri yön limitini yaklaşık $n(t)\cdot0.007$ kW'a kadar aşardı.

\subsection{Dışsal kapasite payı (A6 varsayımı)}\label{app:a6}

EV2Gym, esnek olmayan yükün tepe değerini $m_L\sim\mathrm{Normal}(0.4,0.1^2)$ olmak üzere $m_LP_{\sm{tr}}$'ye ölçekler ve yükü $P_{\sm{tr}}$'de kırpar.
PV üretimi, normalize edilmiş bir $\mathrm{pv}(t)\le0.853\cdot1.1<0.94$ profili ve $m_P\sim\mathrm{Normal}(0.4,0.1^2)$ ile $\mathrm{pv}(t)\,m_PP_{\sm{tr}}$'dir; her ikisi de epizot başına bir kez çekilir.
Dolayısıyla $C(t)\ge P_{\sm{tr}}(1-m_L)$'dir; $P_{\sm{tr}}=\rho_{\sm{tr}}N$ ve $\rho_{\sm{tr}}\ge3$ kW ile $C(t)\ge N\varepsilon_q$ koşulu yalnızca $m_L>1-\varepsilon_q/\rho_{\sm{tr}}\ge0.997$ ise, yani yaklaşık altı standart sapmalık bir olayda sağlanamaz.
Benzer şekilde $\Crev(t)\ge P_{\sm{tr}}(1-0.94\,m_P)$ yalnızca $m_P>1.06$ ise, yani altı standart sapmadan büyük bir olayda negatif olur.
Dolayısıyla A6, epizot başına $2\times10^{-9}$'un altında bir olasılıkla sağlanamaz.
A6 dışsal olduğundan herhangi bir ihlal, yedek kuralın (i), (ii) ya da (iv) durumuna ait bir adım olarak görünür; dolayısıyla sonuç tablolarındaki yedek kural sayıları bunu her tekrarda denetler.

\subsection{Senaryolar, varış profilleri ve gözlem}\label{app:scenarios}

\Tab{tab:scenarios} üzerinde görüldüğü gibi şebeke büyüklükleri ve varış hızı şarj ünitesi sayısı $N$ ile doğrusal olarak ölçeklenir, fiyatlar ise ölçeklenmez.
EV2Gym'in üç varış profili varış zamanları, kalış süreleri ve talep edilen enerjiler bakımından farklılaşır.

İşyeri varışları hafta içi sabahlarında yoğunlaşır, kamusal varışlar öğleden sonra ve akşam tepe noktasıyla gün boyuna yayılır, özel varışlar ise geç öğleden sonra ve akşam gerçekleşir ve gece boyu kalışları içerir.
EV2Gym, ayrılışı epizot sonuna dört adımdan daha yakın düşecek bir EV üretmez; dolayısıyla epizot çok geç başladığında gece boyu süren oturumlar sessizce eksik kalır.
Bu nedenle kamusal ve özel profiller 05:00 yerine 12:00'de başlar.
Yalnızca senaryo üreteci üzerinde, düşürülen üretim denemelerinin payı bu durumda \%6.8 ve \%3.7'ye düşer; profil başına üretim çarpanları ortalama şarj ünitesi doluluğunu yaklaşık 0.21 düzeyinde eşitler.
Bu çarpanlar şarj ünitesi başına enerji talebini de yaklaşık \%5 içinde eşitler; bu senaryo üreteci istatistikleri sonuç değil tasarım denetimleridir \TODO{yazarlar bunları yeniden ölçüp doğrular}.
Epizot boş başlar; yükler ve PV her epizot için yeniden çekilir.
Literatür yapılandırmamız, EV belirtim dosyamızı kullanarak EV2Gym'in 100 kW'lık, işyeri profilli, yük ve PV içermeyen 25 şarj üniteli V2G kâr en büyükleme durumunu yeniden üretir.

\begin{table}[t]
\centering
\caption{Senaryo parametreleri. ``Doğrusal'' olarak işaretlenen büyüklükler şarj ünitesi sayısı $N$ ile ölçeklenir.}
\label{tab:scenarios}
\footnotesize
\setlength{\tabcolsep}{4pt}
\begin{tabular}{@{}>{\raggedright}p{3.4cm}>{\raggedright}p{9.6cm}l@{}}
\toprule
Öge & Ayar & Ölçekleme \\
\midrule
Simülatör & EV2Gym, commit \texttt{55ef368c}, değiştirilmemiş; şarj ünitesi başına bir EV & -- \\
İstasyon boyutları & eğitim $N\in\{10,20\}$; test $N\in\{5,10,20,50,100,200\}$ & -- \\
Trafo anma gücü & $P_{\sm{tr}}=\rho_{\sm{tr}}N$; $\rho_{\sm{tr}}\in\{4,7\}$ kW (sıkı, orta); $N\in\{20,200\}$ için 3 kW stres düzeyi & doğrusal \\
Esnek olmayan yük & ortalama tepe $0.4P_{\sm{tr}}$, $P_{\sm{tr}}$'de kırpılır, her epizotta yeniden çekilir & doğrusal \\
PV üretimi & ortalama tepe yaklaşık $0.38P_{\sm{tr}}$, her epizotta yeniden çekilir & doğrusal \\
Talep tepkisi & kapalı & -- \\
Adım, ufuk & $\Delta t=15$ dk, $T=112$ (28 h); duyarlılık 5 dk, $T=336$ & -- \\
Profiller (başlangıç, üretim çarpanı, günler) & işyeri (05:00, 5, hafta içi); kamusal (12:00, 4, tüm günler); özel (12:00, 2.5, tüm günler) & şarj ünitesi başına \\
Ortalama doluluk & tüm profiller için yaklaşık 0.21 & -- \\
EV modelleri & V2G destekli 8 model; verim 0.90 (bir model için 0.93); $E^{\min}=5$ kWh; en kısa kalış 180 dk & -- \\
Şarj gücü azalması çarpanı & 5 (EV2Gym varsayılanı) & -- \\
İstenen enerji & tam batarya; duyarlılık \%80 şarj durumu & -- \\
Fiyatlar & Hollanda gün öncesi fiyatları; deşarj fiyatı $=$ şarj fiyatı; 0.9 çarpanı ve 0.05 \texteuro/kWh enerji akışı maliyetiyle duyarlılıklar & -- \\
Tahmin gürültüsü & Gauss, standart sapma değerin \%10'u & -- \\
Test ızgarası & 36 yapılandırma ($N$, $\rho_{\sm{tr}}$, profil), 6 stres yapılandırması ve literatür yapılandırması & -- \\
Test epizotları & $N\le50$ ve literatür yapılandırması için yapılandırma başına 100, $N\ge100$ için 50; politika başına 3550 & -- \\
\bottomrule
\end{tabular}
\end{table}

\Tab{tab:features}, iki belirteç türünün özniteliklerini listelemektedir.
Boş bir şarj ünitesi için tüm öznitelikler 0'dır ve boş şarj üniteleri ayrıca dikkat anahtarı olarak maskelenir.
Tahminler, bağımsız standart normal $\epsilon_{t,k}$ ile $\hat P(t+k)=\max\{0,P(t+k)(1+0.1\epsilon_{t,k})\}$ biçimindedir; epizot sonunun ötesinde son gerçek değer korunur.
$P_{\sm{tr}}^{\sm{nom}}=\max_tP_{\sm{tr}}(t)$ ve $N_{\sm{day}}=24/\Delta t$ olmak üzere her küresel öznitelik $N$ ile büyümeyen bir orandır.
Varış profili bir öznitelik değildir.

\begin{table}[t]
\centering
\caption{Şarj ünitesi belirtecinin ($F_c=12$) ve küresel belirtecin ($F_g=38$) öznitelikleri ve normalizasyonları.}
\label{tab:features}
\footnotesize
\setlength{\tabcolsep}{3pt}
\begin{tabular}{@{}>{\raggedright}p{7.6cm}>{\raggedright\arraybackslash}p{7.6cm}@{}}
\toprule
Şarj ünitesi belirteci $\xi_i$ & Küresel belirteç $\xi_0$ \\
\midrule
1: doluluk bayrağı & 1--2: $c^{\sm{ch}}(t)$, $c^{\sm{dis}}(t)$ $/c^{\sm{ref}}$ \\
2--3: $E_i(t)$, $E_i^{\sm{cap}}$ $/E^{\sm{ref}}$ & 3--10: $c^{\sm{ch}}(t+k)$, $k=1,\dots,8$, $/c^{\sm{ref}}$ \\
4: şarj durumu $E_i/E_i^{\sm{cap}}$ & 11--12: günün saatinin sinüsü ve kosinüsü \\
5: kalan gerekli enerji $(\Ereq_i-E_i)^{+}/E^{\sm{ref}}$ & 13: hafta sonu bayrağı; 14: epizodun kalan oranı $(T-t)/T$ \\
6: kalan adımlar $\tau_i$ (konformal varyant: $\tau_i^{\sm{lo}}$) $/N_{\sm{day}}$ & 15--16: $C(t)$, $\Crev(t)$ $/P_{\sm{tr}}^{\sm{nom}}$ \\
7: esneklik payı $\chi_i$, $[-N_{\sm{day}},2N_{\sm{day}}]$ aralığına kırpılmış, $/N_{\sm{day}}$ & 17--18: ölçülen $P^{\sm{inf}}(t)$, $P^{\sm{pv}}(t)$ $/P_{\sm{tr}}^{\sm{nom}}$ \\
8--9: $\Pnom_i$, $\Pdis_i$ $/P^{\sm{ref}}$ & 19--34: $k=1,\dots,8$ için $P^{\sm{inf}}$ ve $P^{\sm{pv}}$ tahminleri, $/P_{\sm{tr}}^{\sm{nom}}$ \\
10: fiziksel alt limit $\pmn_i(E_i)$ $/P^{\sm{ref}}$ & 35: doluluk oranı $n(t)/N$; 36: şarj ünitesi başına anma gücü $P_{\sm{tr}}^{\sm{nom}}/(NP^{\sm{ref}})$ \\
11--12: sınırlar $l_i$, $u_i$ $/P^{\sm{ref}}$ & 37--38: $\sum_il_i$, $\sum_iu_i$ $/P_{\sm{tr}}^{\sm{nom}}$ \\
\bottomrule
\end{tabular}
\end{table}

% ---------------------------------------------------------------------
\section{Talep işleme hattı ve tam bilgili optimumlar}\label{sec:pipeline}
% ---------------------------------------------------------------------

\subsection{İstasyon düzeyinde talep işleme hattı}\label{app:admission}

Varış sınırlaması ve kabul kontrolü denetleyiciye değil, istasyon işletmecisinin talep ayarına aittir.
Her test tekrarı için bunlar herhangi bir denetleyici çalışmadan önce bir kez hesaplanır ve ortaya çıkan kabul edilen talepler $\Ereq_i$, karşılaştırılan her yöntem ve her iki optimum tarafından paylaşılır.
Sarmalayıcı, kabul edilen gereksinimi varışta EV2Gym'in EV nesnesine yazar; böylece EV-GNN ve Gurobi eMPC gibi harici karşılaştırma yöntemleri de aynı hedefi izler.
HiGHS eMPC, ufku içindeki varışlar için paylaşılan girdileri kullanır.
Konformal geçişler, paylaşılan gereksinimi kendi $\hat d_i^{\sm{lo}}$ sınırlarıyla yeniden sertifikalar; dolayısıyla bu geçişlerin gereksinimi hiçbir zaman daha büyük değildir; kabul kontrolü içermeyen geçişler ise varışta sınırlandırılmış talebi korur.
Her geçişte eksik enerji, paylaşılan kabul edilen gereksinime göre ölçülür.

Varışlar zaman sırasına göre işlenir; eşitlik hâlinde şarj ünitesi indisi belirleyicidir.
Varan bir EV $j$; dışsal varış enerjisi, istenen enerjisi, ayrılış adımı ve EV parametreleriyle girer ve \Eq{M-eq:cap} ile verilen sınırlama, sertifikalı bir eşik tablosuyla birlikte ilk gereksinimini verir.
Ardından kabul kontrolü, varış adımı $t_0$'da, daha önce kabul edilmiş ve hâlâ bağlı olan EV'ler ile EV $j$ üzerinde \Eq{eq:admission} ile verilen olurluluk LP'sini çözer:
\begin{equation}\label{eq:admission}
\begin{aligned}
&E_i(t+1)=E_i(t)+\Delta t\,p_i(t),\qquad E_i^{\min}\le E_i(t+1)\le E_i^{\sm{cap}},\qquad E_i(d_i+1)\ge E_i^{\sm{tgt}},\\
&-\Pdis_i\le p_i(t)\le\Pnom_i,\qquad p_i(t)\le\varrho_i^{\sm{in}}\,(E_i^{\sm{cap}}-E_i(t))\ (t>t_0),\qquad p_i(t_0)\in[\pmn_i,\bar P_i],\\
&-\hat C_{\sm{rev}}(t)\le\textstyle\sum_ip_i(t)\le\hat C(t)-n(t_0)\varepsilon_q .
\end{aligned}
\end{equation}
Önceki EV'lerin $t_0$'daki enerjileri herhangi bir denetleyicinin yörüngesinden değil, nominal bir plandan okunur.
Her kabul kararından sonra aynı LP'nin esnek bir sürümü bu planı, hedeflere olabildiğince erken ulaşılacak şekilde günceller.
$\hat C$ ve $\hat C_{\sm{rev}}$ limitleri $t_0$'daki ölçülen değerleri, sonraki $H_f$ adım için gürültülü tahminleri ve sonrasında son tahmini kullanır.
Eğim $\varrho_i^{\sm{in}}$, tüm $E$ değerleri için $\bar P_i(E)\ge\min\{\Pnom_i,\varrho_i^{\sm{in}}(E_i^{\sm{cap}}-E)\}$ koşulunu sağlayan en büyük değerdir; dolayısıyla doğrusal şarj gücü azalması modeli bir iç yaklaşıklıktır ve LP açısından olurlu her çizelge fiziksel olarak gerçeklenebilir.

LP içinde yeni EV'nin hedefi $E_j^{\sm{tgt}}=\Ereq_j$'dir; önceki bir EV'nin hedefi ise $\Ereq_i$ ile $\Psi_i^{\sm{LP}}(E_i(t_0))-10^{-6}$ değerlerinin küçüğüdür.
Burada $\Psi_i^{\sm{LP}}$, LP'nin şarj gücü azalması modelinde erişilebilen en büyük ayrılış enerjisidir ve bu hedef yalnızca kabul LP'si içinde kullanılır.
Tam batarya talepleri doğrudan hedef olarak alınsaydı, böyle bir EV bağlı olduğu her durumda LP olursuz olurdu; çünkü doğrusal şarj gücü azalması modeli sonlu sayıda adımda hiçbir zaman tam şarja ulaşmaz.
Olurluluk yeni talepte monotondur; dolayısıyla $E_j(d_j+1)$'i en büyükleyen tek bir LP, kabul edilebilir en büyük talebi verir.
Daha önce kabul edilen talepler hiçbir zaman azaltılmaz ve EV $j$'nin düşürülmüş talebi yeniden sınırlama rutininden geçer; dolayısıyla A4 her EV ve her yöntem için sağlanır.
LP en düşük talep için bile olursuz ise talep varış enerjisine eşitlenir.

Her karar yalnızca $t_0$'da mevcut olan bilgiyi kullanır; dolayısıyla bir istasyon işletmecisi işleme hattını çevrim içi olarak çalıştırabilir.
LP olurlu ise, başka bir EV gelmemesi koşuluyla, öngörülen limitler altında nominal plan enerjilerinden başlayarak tüm hedefleri karşılayan bir açık döngü çizelge vardır.
İşleme hattı hiçbir denetleyici için A5'i garanti etmez; çünkü gerçekleşen enerjiler nominal planın altında kalabilir, tahminler gürültülüdür ve sonraki varışlar yük ekler.
Bu nedenle Teorem~\ref{M-thm:invariance} kapsamına girmez; etkisini kabul kontrolü içermeyen değerlendirme geçişleri ölçer.

\subsection{Optimumların doğrusal programları}\label{app:oracle}

Her iki tam bilgili LP de epizodun tüm oturumlarını varış ve ayrılış adımları, varış enerjileri ve kabul edilen gereksinimleriyle; tüm fiyatları ve gerçekleşen $C(t)$ ve $\Crev(t)$ limitlerini birlikte bilir.
Her $i$ oturumu ve $t\in\mathcal T_i=\{t_i^{\sm{arr}},\dots,d_i\}$ adımı için değişkenler $p_i^{+}(t),p_i^{-}(t)\ge0$ ve $E_i(t)$'dir; bunlara bir eksik enerji gevşek değişkeni $z_i\ge0$ ve yumuşak LP'de aşırı yük gevşek değişkenleri $o^{\pm}(t)\ge0$ eşlik eder.
$G_i^{\max}=0.01\,(d_i-t_i^{\sm{arr}}+1)$ kWh yuvarlama payıyla kısıtlar şunlardır:
\begin{equation}\label{eq:lp}
\begin{aligned}
&E_i(t_i^{\sm{arr}})=E_i^{\sm{arr}},\qquad E_i(t+1)=E_i(t)+\Delta t\,\big(p_i^{+}(t)-p_i^{-}(t)\big),\qquad E_i^{\min}-G_i^{\max}\le E_i(t+1)\le E_i^{\sm{cap}},\\
&p_i^{+}(t)\le\Pnom_i,\qquad p_i^{-}(t)\le\Pdis_i,\qquad p_i^{+}(t_i^{\sm{arr}})\le\bar P_i(E_i^{\sm{arr}}),\qquad p_i^{+}(t)\le\varrho_i^{\sm{out}}\big(E_i^{\sm{cap}}-E_i(t)\big)\ (t>t_i^{\sm{arr}}),\\
&E_i(d_i+1)+z_i\ge\Ereq_i-G_i^{\max},\qquad -\Crev(t)-o^{-}(t)\le\textstyle\sum_{i}\big(p_i^{+}(t)-p_i^{-}(t)\big)\le C(t)+o^{+}(t),
\end{aligned}
\end{equation}
burada katı LP $o^{\pm}=0$ olarak sabitler.
Bir çizelgenin kârı $\mathrm{Pr}(p)=\sum_t\Delta t\sum_i\big(c^{\sm{dis}}p_i^{-}-c^{\sm{ch}}p_i^{+}-c^{\sm{thr}}(p_i^{+}+p_i^{-})\big)$'dir; enerji akışı maliyeti $c^{\sm{thr}}$ tek bir duyarlılık çalışması dışında sıfırdır.

\begin{lemma}[Şarj gücü azalması satırları]\label{lem:taperrows}
Her EV, Lemma~\ref{lem:a3} koşullarını sağlasın ve $\varrho_{\kappa,i}=(1-e^{-\kappa_i})/\Delta t$ olsun.
Bu durumda her $E\in[E^{\min},E^{\sm{cap}}]$ için $\bar P_i(E)\le\min\{\Pnom_i,\varrho_{\kappa,i}(E_i^{\sm{cap}}-E)\}$'dir.
Eşitlik, tek adımlık geçiş bandı $\varsigma\in(\varsigma_i^{\sm{tr}}-\delta_i,\varsigma_i^{\sm{tr}})$ dışında sağlanır ve $\sup_E\bar P_i(E)/(E_i^{\sm{cap}}-E)=\varrho_{\kappa,i}$'dir.
\end{lemma}

\begin{proof}
Şarj durumu birimlerinde \Eq{eq:taper} gereğince $\bar P\Delta t/E^{\sm{cap}}=\varphi/E^{\sm{cap}}-\varsigma$'dır ve $\varsigma\ge\varsigma^{\sm{tr}}$ için bu kazanç $\min\{\delta,(1-\varsigma)(1-e^{-\kappa})\}$ olup sağ tarafa eşittir.
$\varsigma+\delta\le\varsigma^{\sm{tr}}$ sabit akım bölgesinde kazanç $\delta$'dır ve sağ taraf ancak ve ancak $\delta\le(1-\varsigma)(1-e^{-\kappa})$ ise $\delta$'ya eşittir.
$1-\varsigma\ge1-\varsigma^{\sm{tr}}=m\delta/\kappa$ olduğundan $\kappa\le m(1-e^{-\kappa})$ olması yeterlidir; $\kappa=mz$ ile bu $z\le1-e^{-mz}$, yani $\kappa\le\kappa^{*}(m)$ demektir.
Geçiş bandında $y=\varsigma^{\sm{tr}}-\varsigma\in(0,\delta)$ olsun; böylece $1-\varphi/E^{\sm{cap}}\ge(1-\varsigma^{\sm{tr}})e^{-m(\delta-y)/(1-\varsigma^{\sm{tr}})}$ olur.
$e^{v}\ge1+v$ ve $m\ge1$ ile bu sınırın sağ tarafı en az $e^{-\kappa}(1-\varsigma^{\sm{tr}}+my)\ge e^{-\kappa}(1-\varsigma)$'dır.
Dolayısıyla kazanç en fazla $(1-\varsigma)(1-e^{-\kappa})$'dır ve zaten en fazla $\delta$'dır; dolayısıyla eşitsizlik her yerde sağlanır ve supremum üst bölgedeki eşitlikten çıkar.
\end{proof}

Kod, kesin değerden \%1 daha gevşek, geçerli bir dış satır olan $\varrho_i^{\sm{out}}=1.01\,\varrho_{\kappa,i}$'yi kullanır.
Kodla birlikte yayımlanan bir denetim betiği, 320 yapılandırmada eşitsizliği $1.7\times10^{-13}$ kW duyarlılığında ve bant dışındaki eşitliği doğrulamaktadır.
Bu satırlarla LP'ler, EV2Gym'in şarj limitini geçiş bandının adım ayrıklaştırması ve \%1'lik pay dışında modeller.
Kalan gevşetmeler yuvarlama payı ve eşzamanlı şarj ile deşarjdır; ikincisi, ana ızgaradaki gibi $\eta=1$, $c^{\sm{dis}}=c^{\sm{ch}}$ ve $c^{\sm{thr}}=0$ olduğunda maliyet ve enerji açısından nötrdür ve dolayısıyla kesindir.
Enerji akışı maliyeti içeren ya da negatif fiyatlarda 0.9 deşarj fiyatı çarpanı kullanan duyarlılık çalışmalarında LP'ler geçerli üst sınırlar olarak kalır.
Bu durumda eşzamanlı şarj ve deşarj kullanabilirler ve bu tür adımların sayısı raporlanır.

\subsection{Üst sınırlar}\label{app:bounds}

Katı LP, HiGHS \cite{huangfu2018parallelizing} ile sözlükbilimsel olarak çözülür.
Aşama 1, $G^{\max}=0$ ile en küçük toplam eksik enerjiyi $z^{*}=\min\sum_iz_i$ hesaplar; aşama 2 ise yuvarlama payı ve $\sum_iz_i\le z^{*}+10^{-6}$ ile $\mathrm{Pr}(p)$'yi en büyükleyerek $J^{\sm{LP,ceil}}$'i verir.
$G^{\max}=0$ ile aynı aşama 2, yalnızca \Sec{sec:results-supp} içinde raporlanan $J^{\sm{LP}}$'yi verir.
Buna karşılık yumuşak LP, \Eq{eq:lp} altında $\mathrm{Pr}(p)-w_E\sum_iz_i-w_P\Delta t\sum_t(o^{+}(t)+o^{-}(t))$ ifadesini tek aşamada en büyükleyerek $J^{\sm{pen,LP}}$'yi verir.

\begin{proposition}[Katı optimum]\label{prop:oracle}
Lemma~\ref{lem:a3} yapılandırması geçerli olsun ve $\pi$, raporlanan güçleri her adımda \Eq{M-eq:transformer} koşulunu sağlayan ve kabul edilen gereksinimlere göre gerçekleşen toplam eksik enerjisi en fazla $z^{*}$ olan herhangi bir denetleyici olsun.
Bu durumda $J^{\pi}\le J^{\sm{LP,ceil}}$'dir.
\end{proposition}

\begin{proof}
$\tilde p_i(t)$ raporlanan güçler ve $E_i^{\sm{nom}}(t)=E_i^{\sm{arr}}+\Delta t\sum_{s<t}\tilde p_i(s)$ yalnızca faturalandırılan güçlerin üreteceği enerji olsun.
EV2Gym tam olarak $\tilde p$'yi faturalandırır; dolayısıyla $J^{\pi}$, $p^{\pm}=\tilde p^{\pm}$ noktasındaki $\mathrm{Pr}$'ye eşittir.
Yuvarlama etkin adım başına 0.01 kWh'ten daha az ekler; dolayısıyla yörünge boyunca $E_i^{\sm{real}}(t)-G_i^{\max}\le E_i^{\sm{nom}}(t)\le E_i^{\sm{real}}(t)\le E_i^{\sm{cap}}$'dir.
Bu nedenle düşürülmüş enerji tabanı ve düşürülmüş ayrılış satırı, $z_i=(\Ereq_i-E_i^{\sm{real}}(d_i+1))^{+}$ ile $E^{\sm{nom}}$ için sağlanır; güç satırları da $\tilde p_i\le\bar P_i(E_i^{\sm{real}})\le\Pnom_i$ ve $\tilde p_i\ge-\Pdis_i$ olduğundan sağlanır.
Ayrıca $E^{\sm{nom}}(t_i^{\sm{arr}})=E_i^{\sm{arr}}$ olduğundan ilk adım satırı sağlanır.
Şarj gücü azalması satırları için EV2Gym hiçbir zaman $\varphi(E)-E$'den fazlasını teslim etmez; dolayısıyla Lemma~\ref{lem:taperrows} $\tilde p_i(t)\le\varrho_{\kappa,i}(E^{\sm{cap}}-E_i^{\sm{real}}(t))\le\varrho_i^{\sm{out}}(E^{\sm{cap}}-E_i^{\sm{nom}}(t))$ sonucunu verir.
Trafo satırları varsayımın kendisidir ve $\sum_iz_i\le z^{*}$'dır; dolayısıyla yörünge aşama 2 için olurludur ve amaç değeri $J^{\pi}$'ye eşittir.
\end{proof}

\begin{proposition}[Yumuşak cezalı optimum]\label{prop:softoracle}
Lemma~\ref{lem:a3} yapılandırması geçerli olsun ve komutlar \Eq{M-eq:limits} ile verilen fiziksel limitlere kırpılsın.
Bir tekrar üzerinde simüle edilen her yörünge için, eksik enerji paylaşılan kabul edilen gereksinime göre ölçüldüğünde, $J^{\sm{pen}}(\pi)\le J^{\sm{pen,LP}}$'dir.
\end{proposition}

\begin{proof}
$p^{\pm}=\tilde p^{\pm}$, $E=E^{\sm{nom}}$, $z_i=\Delta E_i^{\sm{sh}}$ alınsın ve $o^{\pm}(t)$ her iki yöndeki gerçekleşen aşırı yüklere eşit olsun.
Enerji, güç, ilk adım ve şarj gücü azalması satırları Önerme~\ref{prop:oracle} ispatındaki gibi sağlanır; çünkü bunlar denetleyicinin uyduğu herhangi bir kısıtı değil, yalnızca EV2Gym'in fiziğini kullanır.
Ayrılış satırı için $E^{\sm{nom}}(d_i+1)+z_i\ge E^{\sm{real}}(d_i+1)-G_i^{\max}+z_i\ge\Ereq_i-G_i^{\max}$'dir; trafo satırları ise $o^{\pm}$ tanımı gereği sağlanır.
Bu noktada amaç değeri, aynı faturalandırılan güçler, eksik enerji ve aşırı yük enerjisiyle $J^{\sm{pen}}(\pi)$'ye eşittir.
\end{proof}

Sınır denetleyicinin değil ölçümün bir özelliğidir; dolayısıyla eksik enerji ya da aşırı yük içeren denetleyicileri, konformal ve kabul kontrolsüz geçişleri ve harici karşılaştırma yöntemlerini kapsar.
Sarmalayıcıdaki her geçiş fiziksel kırpmayı uygular ve Gurobi eMPC bağdaştırıcısı da komutlarını aynı şekilde kırpar.
Kodla birlikte yayımlanan bir birim testi, rastgele eylemler de dâhil olmak üzere küçük tekrarlarda tüm kural tabanlı yöntemler için $J^{\sm{pen}}\le J^{\sm{pen,LP}}$ koşulunu denetler.
EV2Gym içinde bir denetleyici faturalandırılmayan yuvarlama kazancı yoluyla $J^{\sm{LP}}$'yi aşabilir; bu nedenle ana metin garantinin maliyeti için $J^{\sm{LP,ceil}}$'i kullanır ve $J^{\sm{LP}}$'yi yalnızca burada raporlar.

% ---------------------------------------------------------------------
\section{Karşılaştırılan yöntemler, eğitim ve istatistikler}\label{sec:training-supp}
% ---------------------------------------------------------------------

\subsection{Karşılaştırılan yöntemlerin gerçeklenmesi}\label{app:baselines}

Her yöntem ve her değerlendirme geçişi \Sec{sec:pipeline} içinde tanımlanan paylaşılan kabul edilen gereksinimi alır ve eksik enerji her zaman bu gereksinime göre ölçülür.
Güvenlik katmanı içermeyen yöntemler yalnızca fiziksel kırpma $\clip(x_i,\pmn_i,u_i)$ işlemini uygular; katman içeren varyantlar ise çıktılarını ham eylem $x$ olarak \Alg{alg:safety} üzerinden geçirir ve böylece Teorem~\ref{M-thm:invariance} sonucunu devralır.

\textbf{Cezaya dayalı yöntemler.}
Attention+penalty, MLP+penalty ve ayrılış sınırı içermeyen varyant, $r(t)$ ödülünden iki ceza çıkarılarak elde edilen değerden öğrenir.
Bu cezalar, $t$ adımında ayrılan EV'lerin ayrılış eksikliğinin $w$ katı ve adımın trafo aşırı yük enerjisinin $w$ katıdır.
Eksik enerji kabul edilen gereksinime göre ölçülür, aşırı yük sıkılaştırılmamış limitleri ve raporlanan güçleri kullanır ve $w\in\{1,10,100\}$\,\texteuro/kWh ağırlığı ayarlanır.
En yüksek saatlik fiyatın on bir katından fazla olan 10\,\texteuro/kWh değerlendirme ağırlığında, cezalı amacın optimumunda hiçbir ihlal kârlı değildir; ancak öğrenilmiş politikalar yine de ihlalde bulunabilir.

Attention+penalty ve MLP+penalty $\beta=0$ kullanır ve aktör kayıpları yalnızca eylemi komut edilen gücü etkileyemeyen belirteçleri maskeler.
Ayrılış sınırı içermeyen varyant güvenlik katmanını $l_i=\pmn_i$ ile çalıştırır; dolayısıyla $\beta=0.01$ değerini korur ve yedek kural adımlarını \method{} gibi maskeler.

\textbf{Attention+PPO-Lagrangian.}
Bu karşılaştırma yöntemi, beklenen epizot eksik enerjisi ve aşırı yük enerjisinin en fazla sıfır olması koşuluyla $\max J_r$ kısıtlı problemini çözer \cite{altman2021constrained,tessler2018reward}.
Dikkat tabanlı politika fiziksel kırpma ve kâr ödülüyle çalışır; kritiğin biri değer, diğer ikisi her bir maliyet için olmak üzere, her biri kendi avantaj kestirimine sahip üç başlığı vardır.
PPO, mini yığın başına normalize edilen $(A_r-\Lambda_EA_E-\Lambda_PA_P)/(1+\Lambda_E+\Lambda_P)$ birleşik avantajını kullanır.
Her veri toplama turunda bir kez, $k\in\{E,P\}$ için $\Lambda_k\gets\clip(\Lambda_k+\eta_\Lambda(\hat J_k-0),0,1000)$ güncellemesi yapılır; başlangıçta $\Lambda_k=0$'dır ve adım büyüklüğü $\eta_\Lambda\in\{0.5,2,8\}$ ayarlanır.
Sıfır maliyet limiti problemin katı gereksinimleriyle örtüşür; 1.000 üst sınırı ihlaller sürdüğü sürece çarpanları sonlu tutar ve eğitim kayıtları bunların yörüngelerini kaydeder.

\textbf{Boyuttan bağımsız ve sabit boyutlu ağlar.}
Deep Sets+safety layer her şarj ünitesi belirtecini ortak bir ağla gömer ve bağlam olarak tüm gömmelerin maskeli ortalamasını ve küresel belirteci ekler.
Belirteç başına bir Gauss dağılımı çözer; dolayısıyla hiçbir belirteç diğerine dikkat etmez.
MLP karşılaştırma yöntemleri, 20 şarj ünitesi belirteci ile küresel belirtecin sıfırla doldurulmuş birleşimini 640 birimlik iki gizli katmandan geçirir.
Bu ağların genişliği, dikkat tabanlı aktörün parametre sayısıyla eşleşmektedir (614.440'a karşı 610.370); MLP karşılaştırma yöntemleri yalnızca $N\le20$ için değerlendirilir ve raporlanır.


\textbf{Kurallar.}
Olabildiğince hızlı (AFAP), olabildiğince geç (ALAP), döngüsel sıralama (round robin) ve LLF kuralları fiyattan habersizdir ve yalnızca şarj eder.
sLLF \cite{chen2022smoothed}, her EV için sonraki adım esneklik payını $\chi_i^0=(E_i-\Bh_i(\tau_i-1))/(\Pnom_i\Delta t)$ ve $p_i=\clip(\Pnom_i(L-\chi_i^0),0,u_i)$ hızını hesaplar; burada ortak düzey $L$, $\sum_ip_i=\max(0,\Ctil)$ olacak şekilde ikiye bölme yöntemiyle bulunur.
Bu, tepe hızı $\Pnom_i$ ve şarj gücü azalmasını içeren bir üst sınırla birlikte yayımlanmış kuraldır.
Fiyat eşiği kuralı mevcut fiyatı, o ana kadar görülen tüm fiyatların ve bilinen fiyat penceresinin medyanıyla karşılaştırır.
Medyanın altında $\Ctil(t)$ değerini LLF ile paylaştırır, medyanın üstünde en yüksek V2G deşarjını talep eder, aksi hâlde boşta bekler.

\textbf{Model öngörülü karşılaştırma yöntemleri.}
eMPC-HiGHS, EV2Gym'in eMPC'sinin \cite{diazlondono2026open} HiGHS ile kayan ufuklu bir net güç LP'si olarak yaptığımız yeniden gerçeklemesidir.
LP'si araç başına limitleri, şarj gücü azalmasını içeren kesin ilk adım sınırlarını ve kabul edilen gereksinimleri 1.000\,\texteuro/kWh'lik esnek bir ayrılış kısıtıyla kullanır.
Fiyatlar ve varışlar 10 adımlık ufku içinde kesin olarak bilinir; yük ve PV ise EV2Gym'in gürültülü tahminlerinden gelir.
24 saatlik ufuklu bir varyant ve Gurobi tabanlı özgün eMPC planlanmış ancak çalıştırılmamıştır.

\textbf{Harici öğrenen karşılaştırma yöntemleri.}
Heterojen çizge politikası (EV-GNN) \cite{orfanoudakis2025scalable}, \texttt{01aa1b7a} commit'indeki açık kaynak kodundan kendi eğitim betiğiyle çalıştırılacak harici, boyuttan bağımsız bir karşılaştırma yöntemi olarak planlanmış, ancak bu çalışmaya dâhil edilmemiştir.
Ayrıca çizge ile güçlendirilmiş karar Transformer'ı (GNN-DT) \cite{orfanoudakis2026graph}, açık kaynak lisansı altında herkese açık koda sahiptir (DT4EVs, commit \texttt{8b327d12}).
GNN-DT, değiştirilmiş bir simülatör kopyası içinde bir optimizasyon çözücüsünün her senaryo için ürettiği yörüngeler üzerinde çevrim dışı eğitilir ve bu iş akışı senaryolarımıza uyarlanmamıştır.

\subsection{Eğitim döngüsü}\label{app:ppo}

\Alg{alg:ppo}, \Alg{alg:safety} içindeki güvenlik katmanının ortamın parçası olarak çalıştığı \method{} eğitim döngüsünü listelemektedir.
Her ortam kendi arabelleğine yazar ve güncelleme anında tüm belirteç tensörleri, yığındaki en büyük $N$'ye kadar maskeli girdilerle sağdan doldurulur.
Doldurma ne gerçek belirteçlerin aktör çıktılarını ne de değeri değiştirir (Lemma~\ref{lem:equivariance}); log-olasılık, entropi ve düzenlileştirici toplamları yalnızca maskeler üzerinden alınır.
Öğrenme ödülü, farklı $N$ değerlerine sahip istasyonların getirileri aynı ölçekte olsun diye $N r^{\sm{ref}}$'ye bölünmüş kârdır; raporlanan her sonuç ise bölünmemiş kârı kullanır.

\begin{algorithm}[t]
\caption{Değişken sayıda şarj ünitesiyle \method{} için PPO eğitimi}\label{alg:ppo}
\begin{algorithmic}[1]
\Require 16 ortam ($N=10$ ile 8, $N=20$ ile 8); aktör $\pi_\theta$ ve kritik $V_\omega$
\While{ortam adımı sayısı $3\times10^{6}$'dan az}
  \For{her ortamda 128 adım}
    \State belirteçleri, $l$, $u$, $\Ctil$, $\Crev$ ve olurluluk bayrağını gözlemle
    \State $a\sim\pi_\theta(\cdot\mid s)$ örnekle; $x\gets P^{\sm{ref}}a$
    \State olurlu ise $p\gets\Pi_t(x)$, değilse yedek kural \Comment{\Alg{alg:safety}}
    \State simülatörü bir adım ilerlet; $r\gets$ adımın gerçekleşen kârı (\Eq{M-eq:objective}); $u_i-l_i>10^{-9}$ olan belirteçler üzerinden $r/(N r^{\sm{ref}})$, $a$, $\log\pi_\theta(a\mid s)$ değerlerini sakla
  \EndFor
  \State avantajları GAE ile hesapla ($\gamma=0.99$, $\lambda_{\sm{GAE}}=0.95$); belirteçleri yığındaki en büyük $N$'ye kadar doldur
  \For{10 dönem, $n_{\sm{mb}}$ geçişlik mini yığınlar}
    \State normalize edilmiş avantajlarla olurlu geçişler üzerinde kırpılmış vekil amaç ($\epsilon_{\sm{clip}}=0.2$), $n_{\sm{mb}}$'ye bölünmüş
    \State kayıp $\gets-$vekil amaç $+\,c_{\sm{vf}}\cdot$değer kaybı $-\,c_{\sm{ent}}\cdot$entropi $+\,\mathcal L_{\sm{reg}}$ (\Eq{M-eq:reg}, yalnızca olurlu adımlar)
    \State gradyan normu 0,5'te kırpılarak Adam adımı; yaklaşık KL ıraksaması 0,03'ü aşarsa dönemleri erken durdur
  \EndFor
\EndWhile
\end{algorithmic}
\end{algorithm}

\subsection{Hiperparametreler ve ayarlama}\label{app:hyper}

\Tab{tab:hyper} hiperparametreleri listelemektedir.
Ayarlanan yedi öğrenen yöntemin (\method{}, Attention+penalty, Attention+PPO-Lagrangian, ayrılış sınırı olmayan Attention+safety layer, Deep Sets+\allowbreak safety layer, MLP+\allowbreak safety layer ve MLP+\allowbreak penalty) her biri 12 yapılandırma almıştır.
Bunlar varsayılan nokta ile, yöntem başına sabit bir örnekleme tohumuyla \Tab{tab:hyper} ızgarasından yerine koymadan örneklenen 11 noktadan oluşur.
Her yapılandırma eğitim karışımı üzerinde 0 tohumuyla $0.5\times10^{6}$ adım eğitilir ve tohumları test tekrarlarından ayrık olan, eğitim yapılandırmaları arasında dönüşümlü olarak ilerleyen 24 doğrulama epizodunda değerlendirilir.
Seçim, her yöntem için $w=10$\,\texteuro/kWh ile hesaplanan ortalama doğrulama $J^{\sm{pen}}$ değerini kullanır.
Konformal varyant ve $N=50$ için eğitilen referans \method{} seçimini devralır; $\beta$ ve mimari, güvenlik katmanı içeren tüm yöntemler için sabittir.
Ayarlama 84 eğitim ve 84 doğrulama işi çalıştırmıştır; seçilen yapılandırmalar kodla birlikte yayımlanmaktadır.

EV2Gym bazı eğitim epizotlarını oluşturamamıştır, çünkü negatif konut yükleri tahmin gürültüsü için negatif bir ölçek vermiştir.
Bu nedenle eğitim tohumları ayarlama işi başına bir kez, eğitim işi başına iki kez yeniden çekilmiş; test tekrarları ise değiştirilmeden oluşturulmuştur.

\begin{table}[t]
\centering
\caption{Hiperparametreler ve ayarlama ızgarası.}
\label{tab:hyper}
\footnotesize
\setlength{\tabcolsep}{3pt}
\begin{tabular}{@{}l>{\raggedright}p{6.2cm}>{\raggedright\arraybackslash}p{6.4cm}@{}}
\toprule
Grup & Parametre & Değer \\
\midrule
Ağ & genişlik $d_m$, katmanlar $n_L$, başlıklar $n_H$ & 128, 3 (ön-LayerNorm), 4 \\
 & ileri besleme genişliği, dropout & 512 (GELU), 0 \\
 & aktör başlığı (ortak), kritik çıkış katmanı & $128\to64\to2$; tek sorgulu dikkat havuzlama, $128\to64\to1$ \\
 & $\log\sigma$ sınırlaması ve başlangıç değeri & $[-5,0.5]$, $-0.5$ \\
 & parametreler (aktör / kritik) & 610.370 / 676.481; küçük ağ toplam 228.067 (başlık genişliği 32) \\
 & duyarlık & float32 ağlar, float64 güvenlik katmanı \\
PPO & optimize edici & Adam, $\epsilon=10^{-5}$, doğrusal öğrenme oranı azaltımı \\
 & öğrenme oranı (ayarlanır) & $\{1,3,10\}\times10^{-4}$ \\
 & mini yığın $n_{\sm{mb}}$ (ayarlanır) & $\{256,512,1024\}$ \\
 & entropi katsayısı $c_{\sm{ent}}$ (ayarlanır) & $\{0,10^{-3},10^{-2}\}$ \\
 & indirim $\gamma$, GAE $\lambda_{\sm{GAE}}$ & 0,99; 0,95 \\
 & kırpma aralığı $\epsilon_{\sm{clip}}$, değer katsayısı $c_{\sm{vf}}$ & 0,2; 0,5 \\
 & ortamlar, veri toplama turu uzunluğu & 16 ($N=10$ ile 8, $N=20$ ile 8), 128 \\
 & dönemler, hedef KL, gradyan normu kırpması & 10; 0,03; 0,5 \\
 & avantaj normalizasyonu & mini yığın başına \\
 & çalıştırma başına adım, tohumlar & $3\times10^{6}$; 2 (tohum 0 ve 1) \\
Ceza içeren ödül & eğitim ağırlığı $w$ (ayarlanır), $\beta$ & $\{1,10,100\}$\,\texteuro/kWh; $\beta=0$, ayrılış sınırı olmayan varyantta 0,01 \\
Lagrange & dual adım $\eta_\Lambda$ (ayarlanır), maliyet limiti, $\Lambda_{\max}$ & $\{0.5,2,8\}$; 0; 1.000 \\
Güvenlik katmanı & $P^{\sm{ref}}$, $\beta$ & 11 kW; 0,01 \\
 & $\varepsilon_q$, $\varepsilon_{\sm{sat}}$, tolerans & 0,008 kW; 0,001 kW; $10^{-9}$ kW ya da kWh \\
Öznitelikler & $E^{\sm{ref}}$, $c^{\sm{ref}}$, $H_f$ & 100 kWh; 0,25 \texteuro/kWh; 8 adım \\
Değerlendirme & $J^{\sm{pen}}$ içindeki $w_E$, $w_P$ ağırlıkları & her biri 10 \texteuro/kWh, tüm yöntemler \\
 & test eylemi & deterministik, $a=\mu_\theta$ \\
Konformal & $n_{\sm{cal}}$, $\alpha$ & yapılandırma başına 5.000; $\{0.01,0.05,0.1,0.2\}$ \\
 & kantil tahmincisi & gradyan artırma, derinliği 3 olan 200 ağaç, 2.000 eğitim oturumu \\
\bottomrule
\end{tabular}
\end{table}

\subsection{İstatistiksel analiz}\label{app:stats}

Öğrenen her yöntem iki tohumla eğitilmiştir; $N=50$ için eğitilen referans da iki tohumla eğitilmiştir.
İki tohumla tohum düzeyindeki bir $t$-testinin serbestlik derecesi birdir; dolayısıyla tohum düzeyinde hiçbir test ya da aralık hesaplanmamaktadır.
Her test tekrarı için öğrenen bir yöntemin değeri, iki tohumu üzerinden ortalamadır (tohum ortalamalı politika) ve tüm yöntemler aynı sabit tekrarlar üzerinde çalışır.
Tablolar, tohum ortalamalı politikanın tekrar ortalamalarının yapılandırmalar üzerinden ortalamasını ve ardından aynı istatistiğin iki tohum üzerindeki aralığını raporlamaktadır.
Deterministik kurallar ve eMPC-HiGHS bir kez çalışır ve aralık taşımaz.
Güven aralıkları, test tekrarlarını her yapılandırma içinde $10^4$ kez yeniden örnekleyen yüzdelik bir bootstrap ile elde edilir; epizot başına açığın çeyrekler arası ortalaması (IQM) $2\times10^3$ yeniden örnekleme kullanır \cite{agarwal2021deep}.

\method{} ile karşılaştırmalar her istasyon boyutu için her sonuç tablosunun ölçütleri üzerinde yapılır; güvenlik katmanı içeren her yöntem ayrıca cezaya dayalı ikiziyle karşılaştırılır.
Birincil test, etki büyüklüğü olarak eşleştirilmiş çiftler sıra-çift serili korelasyonuyla birlikte, tohum ortalamalı politikaların tekrar düzeyindeki farkları üzerinde iki yönlü Wilcoxon işaretli sıralar testidir \cite{wilcoxon1945individual}.
Holm düzeltmesi bir test dosyasındaki tüm karşılaştırmalar üzerinden uygulanır ve her dosya, aile başına en fazla 351 karşılaştırmayla bir sonuç tablosunu kapsar.
Wilcoxon testi farkların medyan kaymasıyla, bootstrap aralığı ise ortalamalarıyla ilgilidir.
Büyük cezalı birkaç epizot medyanı kaydırmadan ortalamayı kaydırabilir; dolayısıyla ikisi uyuşmayabilir ve bu tür karşılaştırmalar sonuçsuz olarak raporlanır.
Tohumlar arası değişkenlik hiçbir test tarafından yakalanmaz ve az sayıda tohuma dayanan sonuçlar yanıltıcı olabilir \cite{henderson2018deep}; bu nedenle tohum aralıkları ve tohum sayısı bir sınırlılık olarak raporlanmaktadır.

\subsection{Hesaplama}\label{app:compute}

Çalıştırmalarda 8 çekirdekli ve 16 iş parçacıklı bir AMD Ryzen 7 7700 ile 16 GB'lık bir NVIDIA GeForce RTX 5070 Ti kullanılmıştır.
WSL2'ye 12 iş parçacığı ve 31 GB bellek atanmıştır.
Yazılım olarak WSL2 Ubuntu 24.04 ile Windows 11 Pro, Python 3.13.15 ve CUDA 13.0 ile PyTorch 2.14.1 kullanılmıştır.
Senaryoların, kabul tablolarının ve konformal kalibrasyonların hazırlanması 1,6 saat, ayarlama 7,9 saat ve 18 çalıştırmanın eğitimi 10,2 saat sürmüştür.
298 değerlendirme ve 43 optimum işiyle değerlendirme 2,5 saat, analiz ise 6 dakika sürmüştür.
Temel küme 30 yöntem için 153.100 epizot üretmiş ve hiçbir iş başarısız olmamıştır.
\Tab{tab:t6}, çalıştırma başına eğitim süresini ve kontrol adımı başına çıkarım süresini vermektedir.

% ---------------------------------------------------------------------
\section{Ek bulgular}\label{sec:results-supp}
% ---------------------------------------------------------------------
% Tablolar ve şekiller \analysisdir (varsayılan ../../analysis) altından ya da
% tables/ ve figures/ içindeki yerel kopyalardan okunur. Metindeki her sayı
% kaynak hücresiyle birlikte paper/results_number_sources.md dosyasında
% listelenmiştir.

\TODO{Yazarlar raporlanan her değeri analiz çıktılarıyla karşılaştırarak denetler ve bu bölümün ifadelerini gözden geçirir}.

Bu bölümdeki tablolar analiz iş akışı tarafından üretilmektedir.
İş akışının T2, T3, T3b, T3c, T3d, T5, T4, T7, T6, TS1, TS2 ve TS3 dosyaları bu sırayla Tablo~\ref{tab:t2}--\ref{tab:ts3} olarak yer almaktadır.
Değerler tohum ortalamalı politikaların ortalamalarıdır; parantez içinde iki eğitim tohumu üzerindeki aralık verilmiştir ve \dag{} tek bir çalıştırmaya sahip deterministik yöntemleri işaretler.
Örnek gün şekli, planlanan ablasyonlar ve stres anma düzeyi dışındaki duyarlılık çalışmaları çalıştırılmamıştır.

\subsection{Öğrenme eğrileri ve sıfır-atış aktarım}\label{app:res-tables}

\Fig{fig:curves}, öğrenen yöntemlerin eğitim eğrilerini iki tohum üzerindeki bantlarla birlikte göstermektedir.
\Tab{tab:t2}, her istasyon boyutu için \Fig{M-fig:zeroshot} şeklinin dayandığı sıfır-atış sonuçlarını, $N=50$ için eğitilen referansa göre aktarım açığı dâhil olmak üzere listelemektedir.
Bu referans yalnızca $N=50$ için dağılım içidir; dolayısıyla $N=100$ ve $N=200$ için aktarım açığı iki dışdeğerlemeyi karşılaştırır.

\begin{figure}[t]
\centering
\resultfig{figS1_training_curves.pdf}{Öğrenen yöntemlerin eğitim eğrileri.}
\caption{Öğrenen yöntemlerin ortam adımlarına göre eğitim eğrileri; bantlar iki eğitim tohumu üzerindeki aralığı gösterir.}
\label{fig:curves}
\end{figure}

\restable{T2}{\placeholdertable{tab:t2}{İstasyon boyutları $N$ arasında sıfır-atış aktarım}{T2.tex}}

\subsection{Kısıtların sağlanması}\label{app:res-constraints}

Test ızgarası üzerinden birleştirilmiş olarak \Tab{tab:t3}; E4'ü, karşılanmayan enerjiyi, aşırı yükü, \Eq{eq:deficit} ile verilen açık sınırına göre eksik enerjiyi ve sınırlandırılmış ve kabul kontrolüyle düşürülmüş talepleri vermektedir.
Açık sınırı oranı yalnızca tam güvenlik katmanına sahip yöntemleri kapsar, çünkü ayrılış sınırı içermeyen varyantın böyle bir sınırı yoktur.
İstasyon boyutu başına \Tab{tab:t3b}; E4'ü, yedek kural durumlarına göre olursuz adımları ve toplam limitlerin ve ayrılış alt sınırının bağlayıcılık oranlarını vermektedir.
\Tab{tab:t3c} ve \Tab{tab:t3d}, her iki tohumun tüm epizotları üzerinden toplanan olay sayılarını genel olarak ve istasyon boyutu başına vermektedir.
Bu tabloların son sütunu, Teorem~\ref{M-thm:invariance}(b) için epizot düzeyindeki denetimdir: hiç olursuz adım içermeyen epizotlarda gereksinimi karşılanmayan EV'ler.
Bu sütun, ayrılış sınırına sahip konformal olmayan her yöntem için sıfırdır; ayrılış sınırı içermeyen varyantta ise bu tür 1.551 EV bulunmuştur.
Ancak tüm epizotlar üzerinden bu varyant, gereksinimi karşılanmayan daha az EV bırakmıştır: varyant için 1.551, \method{} için 6.158 (\Sec{M-sec:res-bound}).
Konformal varyantlar için denetim, yeniden sertifikalanmış gereksinimi ortak gereksinimin altında kalan ya da öngörülen sınırdan önce ayrılan EV'leri de saymaktadır; teorem bu durumları kapsamaz.

\restable{T3}{\placeholdertable{tab:t3}{Tüm test ızgarası yapılandırmaları üzerinden yöntem başına kısıtların sağlanması}{T3.tex}}
\restable{T3b}{\placeholdertable{tab:t3b}{İstasyon boyutu $N$ başına E4, yedek kural durumlarına göre olursuz adımlar ve bağlayıcılık oranları}{T3b.tex}}
\restable{T3c}{\placeholdertable{tab:t3c}{Değerlendirilen tüm epizotlar üzerinden yöntem başına olay sayıları}{T3c.tex}}
\restable{T3d}{\placeholdertable{tab:t3d}{Yöntem ve istasyon boyutu $N$ başına olay sayıları}{T3d.tex}}

\subsection{Garantinin maliyeti ve kısıt yönetimi}\label{app:res-ablation}

\Tab{tab:t5}, \Fig{M-fig:cost} şeklinin dayandığı garantinin maliyetini dikkat tabanlı ve MLP çiftleri için $J^{\sm{pen}}$ ve kâr cinsinden vermektedir.
\Fig{fig:ablation} ve \Tab{tab:t4}, mimarileri ve kısıt yönetimi varyantlarını eğitim boyutlarında ve sıfır-atış aktarımda karşılaştırmaktadır.
$N\ge50$ için \method{} yönteminin açığı 0,901, Deep Sets+safety layer'ın 0,674 ve Attention+safety layer without departure bound'un 0,770 olmuştur.
Katmanı kullanım aşamasında uygulayan ceza ile eğitilmiş politika 0,910'a, Attention+penalty 0,996'ya ve Attention+PPO-Lagrangian 1,138'e ulaşmıştır.

\restable{T5}{\placeholdertable{tab:t5}{İstasyon boyutu $N$ başına garantinin maliyeti}{T5.tex}}

\begin{figure}[t]
\centering
\resultfig{figS2_ablation.pdf}{Mimari ve kısıt yönetimi varyantlarının açığı.}
\caption{\Tab{tab:t4} içindeki mimari ve kısıt yönetimi varyantlarının eğitim boyutlarında ve sıfır-atış aktarımdaki açığı.}
\label{fig:ablation}
\end{figure}

\restable{T4}{\placeholdertable{tab:t4}{Mimari ve kısıt yönetimi}{T4.tex}}

\subsection{Konformal ayrılış sınırı}\label{app:res-conformal}

\Fig{fig:conformal} ve \Tab{tab:t7}, \method{}-CP yönteminin kapsamasını, ortak gereksinime göre eksik enerjisini, açığını ve şarj edilen enerjisini kapsama dışı kalma düzeyi $\alpha$'ya göre göstermektedir.
0,99, 0,95, 0,90 ve 0,80 hedefleri için elde edilen kapsama 0,992, 0,950, 0,902 ve 0,799 olmuştur.
Tüm epizotlar üzerinden toplandığında eksik enerji $\alpha=0.01$ için 771.853 kWh, $\alpha=0.05$ için 12.507 kWh ve bilinen ayrılış adımlarıyla 2.568 kWh olmuştur.
Karşılaştırma için kalibre edilmemiş nokta tahmini 0,899 kapsamaya ve 1,322 açığına ulaşmıştır.

\begin{figure}[t]
\centering
\resultfig{figS3_conformal_alpha.pdf}{$\alpha$'ya göre kapsama, eksik enerji, açık ve şarj edilen enerji.}
\caption{Konformal ayrılış sınırının ödünleşimi: kapsama dışı kalma düzeyi başına hedefine göre kapsama, ortak gereksinime göre eksik enerji, açık ve şarj edilen enerji.}
\label{fig:conformal}
\end{figure}

\restable{T7}{\placeholdertable{tab:t7}{Konformal ayrılış sınırı}{T7.tex}}

\subsection{Hesaplama süresi}\label{app:res-time}

\Fig{fig:inference} ve \Tab{tab:t6}, kontrol adımı başına çıkarım süresini ve çalıştırma başına eğitim süresini raporlamaktadır.
\method{} yönteminin bir eğitim çalıştırması 1,15 saat sürmüş; bir kontrol adımı $N=10$ için 1,68 ms, $N=200$ için 7,76 ms almıştır.
$N=200$ için Deep Sets+safety layer adım başına 5,99 ms, katmanlı kurallar en fazla 2,72 ms ve eMPC-HiGHS+safety layer 14,13 ms gerektirmiştir.

\begin{figure}[t]
\centering
\resultfig{figS4_inference_time.pdf}{$N$'ye göre kontrol adımı başına çıkarım süresi.}
\caption{Politika ve güvenlik katmanı dâhil, $N$'ye göre kontrol adımı başına çıkarım süresi.}
\label{fig:inference}
\end{figure}

\restable{T6}{\placeholdertable{tab:t6}{Hesaplama}{T6.tex}}

\subsection{Göreli açıklar, stres düzeyi ve literatür yapılandırması}\label{app:res-sensitivity}

\Tab{tab:ts1}, yumuşak cezalı optimum ve her iki katı optimum için göreli açıkları toplamların oranı olarak raporlamaktadır.
$J^{\sm{pen,LP}}$ çoğu zaman sıfıra yakın ya da negatif olduğundan \method{} yönteminin göreli $J^{\sm{pen}}$ açığı $N=5$ ve $N=10$ için \%100'ü aşmıştır; bu nedenle ana metin mutlak açığı kullanmaktadır.
\Tab{tab:ts2}, $N=20$ ve $N=200$ için $\rho_{\sm{tr}}=3$ kW stres düzeyini kapsamaktadır.
Burada \method{} yönteminin E4 değeri \%55,30 ve \%65,91; Deep Sets+safety layer'ınki ise \%65,14 ve \%73,25 olmuştur.
EV2Gym'in 25 şarj üniteli yapılandırmasında Attention+safety layer without departure bound'un açığı 0,363, katmanlı eMPC-HiGHS'nin 0,456 ve \method{} yönteminin 0,507 olmuştur (\Tab{tab:ts3}).

\restable{TS1}{\placeholdertable{tab:ts1}{Toplamların oranı olarak göreli açıklar}{TS1.tex}}
\restable{TS2}{\placeholdertable{tab:ts2}{$\rho_{\sm{tr}}=3$ kW stres düzeyi}{TS2.tex}}
\restable{TS3}{\placeholdertable{tab:ts3}{EV2Gym'in 25 şarj üniteli literatür yapılandırmasındaki sonuçlar}{TS3.tex}}

\subsection{Birleştirilmiş kısıt ölçütleri}\label{app:res-example}

\Fig{fig:constraintspooled}, her yöntem ve trafo anma düzeyi için test ızgarasının tüm istasyon boyutları üzerinden birleştirilmiş karşılanmayan enerjiyi ve aşırı yük enerjisini göstermektedir.
Aşırı yük enerjisi, \Tab{tab:t3c} içinde katmanlı her yöntemin sıfır aşırı yük adımıyla uyumlu olarak yalnızca güvenlik katmanı içermeyen yöntemlerde görülmektedir.
Konformal varyantlar gösterilmemiştir; bunların eksik enerjisini \Tab{tab:t7} vermektedir.

\begin{figure}[t]
\centering
\resultfig{figS6_constraints_pooled.pdf}{İstasyon boyutları üzerinden birleştirilmiş, yöntem başına kısıt ölçütleri.}
\caption{İstasyon boyutları üzerinden birleştirilmiş, yöntem ve trafo anma düzeyi başına kısıt ölçütleri.}
\label{fig:constraintspooled}
\end{figure}


\clearpage
\IfFileExists{IEEEtran.bst}{\bibliographystyle{IEEEtran}}{\bibliographystyle{ieeetr}}
\bibliography{../manuscript_tr/refs}

\end{document}
